
\begin{enumerate}
\def\labelenumi{\arabic{enumi}.}
\setcounter{enumi}{17}
\tightlist
\item
  设 \(G\) 是连通实李群，\(Z\) 是其中心，\(Z_0\) 是 \(Z\) 的单位元连通分支，\(g\) 是 \(G\) 的李代数，而 \(h\) 是 \(g\) 的中心。于是 \(Z\) 和 \(Z_0\) 是 \(G\) 的李拟子群，其李代数为 \(h\)（练习 3）。\(g\) 上的可接受范数，是指定义 \(g\) 的拓扑并使 \(g\) 成为赋范李代数的范数。
\end{enumerate}

\begin{enumerate}
\def\labelenumi{(\alph{enumi})}
\item
  对所有 \(r > 0\) 及 \(z \in \mathfrak{z}\)，存在一个 \(\mathfrak{g}\) 上的可接受范数，其在 \(z\) 处的取值为 \(< r\)。（设 \(q\) 是 \(\mathfrak{g}\) 上的可接受范数。证明对于适当选取的 \(\lambda > 0\)，函数 \(x \mapsto \|x\| = \lambda q(x) + \inf_{y \in \mathfrak{z}} q(x + y)\) 具有所需性质。）
\item
  证明下列条件等价：
\item
  \(Z_{0}\) 是单连通的；
\end{enumerate}

\begin{enumerate}
\def\labelenumi{(\roman{enumi})}
\setcounter{enumi}{1}
\item
  对 g 上的每个可接受范数，\(\exp_{G}\) 在以 0 为中心、半径为 \(\pi\) 的开球上的限制是单射；
\item
  存在 \(r > 0\)，使得对于 \(\mathfrak{g}\) 上的每个可接受范数，\(\exp_{\mathbb{G}}\) 在以 0 为中心、半径为 \(r\) 的开球上的限制是单射。
\end{enumerate}

（为证明 (iii) \(\Rightarrow\) (i)，使用 (a)。为证明 (i) \(\Rightarrow\) (ii)，设 \(x, y\) 属于 \(\mathfrak{g}\)，\(\| x \| < \pi\)，\(\| y \| < \pi\)，\(x \neq y\)，且 \(\exp_{\mathfrak{g}} x = \exp_{\mathfrak{g}} y\)。于是 \(\exp \mathrm{d} x = \exp \mathrm{d} y\)，从而得到 \(\mathrm{ad} x = \mathrm{ad} y\)（将 §6 的命题 17 应用于 \(\mathfrak{g}\) 的复化）。因此存在一个非零 \(z\) 属于 \(\mathfrak{z}\)，使得 \(\exp_{\mathfrak{g}} z = e\)；于是 (i) 为假。）

\begin{enumerate}
\def\labelenumi{(\alph{enumi})}
\setcounter{enumi}{2}
\tightlist
\item
  考察练习 13 (b) 中的群 \(\mathbf{G}'\)，证明若将 \(\pi\) 换成 \(> \pi\) 的数，(b) 的结论不再成立。（在练习 13 的记号下，使用范数 \(ae_{1} + be_{2} + ce_{3} \mapsto (a^{2} + b^{2} + c^{2})^{\frac{1}{2}}\)，定义在 \(\mathrm{L}(\mathbf{G}')\) 上。）
\end{enumerate}

\begin{enumerate}
\def\labelenumi{\arabic{enumi}.}
\setcounter{enumi}{18}
\item
  设 G 是局部紧群，H 是 G 的闭子群。假设 H 是有限维单连通可解李群且 G/H 紧。证明 G 存在一个紧子群 L，使得 G 是 L 与 H 的半直积。（对 dim H 归纳。考虑 H 的最后一个非平凡导出群，并使用《积分理论》第 VII 章 § 3 命题 3。）
\item
  设 \(G\) 是有限维连通可解实李群。引入 §6 第 10 号命题 20 的证明中所用的记号 \(S, S', F, \sigma\)。
\end{enumerate}

\begin{enumerate}
\def\labelenumi{(\alph{enumi})}
\item
  利用命题 21，证明 \(\sigma\) 是从 S 到 S' 的底层实李群的一个李子群的同构。
\item
  推出 \(\tilde{G}\)（即 \(G\) 的万有复化）可识别为 \(S' / \sigma(F)\)，且从 \(G\) 到 \(\tilde{G}\) 的典范映射，是将 \(G\) 同构到 \(\tilde{G}\) 的底层实李群的一个李子群上的映射。
\end{enumerate}

¶ 21. (a) 设 G 是满足下列性质的单连通可解实李群：(α) L(G) 为 n 维，并具有一个 n - 1 维交换理想，该理想对应于 G 的子群 A；(β) G 的中心中存在一个不属于 A 的元素 σ。证明存在元素 \(x\)，属于 \(\mathrm{L}(\mathbf{G})\)，使 \(\exp x = \sigma\)。证明 \(\mathrm{L}(\mathbf{G})\) 是一个交换理想与一个具有如下基的理想的乘积：

\[
(x, a _ {1}, b _ {1}, \dots , a _ {k}, b _ {k})
\]

 其中 \(a_1, b_1, \ldots, a_k, b_k\) 属于 \(\mathrm{L}(\mathrm{A})\)，\([x, a_i] = 2\pi n_i b_i\)，\([x, b_i] = -2\pi n_i a_i\)（\(n_i \in \mathbf{Z} - \{0\}\)）对所有 \(i\) 成立。将练习 13 (d) 的结果推广到 G。

\begin{enumerate}
\def\labelenumi{(\alph{enumi})}
\setcounter{enumi}{1}
\tightlist
\item
  设 G 是有限维单连通可解实李群。设 D 是 G 的中心的离散子群。存在 L(G) 的一个基 \((x_{1}, x_{2}, \ldots, x_{n})\) 及一个整数 \(r \leqslant n\)，具有下列性质：(α) G 的每个元素都能唯一写成如下形式：
\end{enumerate}

\[
\left(\exp t _ {1} x _ {1}\right) \dots \left(\exp t _ {n} x _ {n}\right)
\]

  其中 \(t_1, \ldots, t_n\) 属于 \(\mathbf{R}\)；\((\beta)x_1, \ldots, x_r\) 两两可交换，且

\[
(\exp x _ {1}, \dots , \exp x _ {r})
\]

  是交换群 D 的一个基。（对 G 的维数归纳。设 a 是 L(G) 的极大交换理想，A 是相应的积分子群。于是 A 在 G 中闭，且 DA/A 是 G/A 的中心的离散子群，可以对其应用归纳假设，从而存在 L(G/A) 中的元素 \(x_1^*, \ldots, x_m^*\) 及一个整数 \(s\)。对于 \(1 \leqslant i \leqslant s\)，设 \(\sigma_i\) 是 D 中一个模 A 的类为 exp \(x_1^*\) 的元素。逐个使用 (a)，构造代表元 \(x_1, \ldots, x_m\)，使其成为 \(x_1^*, \ldots, x_m^*\) 的代表元，且 exp \(x_i = \sigma_i\)，以及 \([x_1, x_j] = 0\) 对 \(1 \leqslant i, j \leqslant s\) 成立。

\begin{enumerate}
\def\labelenumi{(\alph{enumi})}
\setcounter{enumi}{2}
\tightlist
\item
  由 (b) 推出每个有限维连通可解实李群都同态于空间 \(\mathbf{R}^n\times \mathbf{T}^m\)（\(m,n\) 为 \(\geqslant 0\) 的整数）。
\end{enumerate}

\begin{enumerate}
\def\labelenumi{\arabic{enumi}.}
\setcounter{enumi}{21}
\item
  设 \(G\) 是有限维连通实李群，且 \(L(G)\) 是约化的。则 \(G = (V \times S)/N\)，其中 \(V\) 是有限维实向量空间，其中 \(S\) 是单连通半单实李群，且 \(N\) 是 \(V \times S\) 的中心离散子群。于是 \(G/\overline{D^1}G\) 同构于 \(V/\mathrm{pr}_1\overline{N}\)。因此 \(G/\overline{D^1}N\) 紧，当且仅当 \(\mathrm{pr}_1\mathrm{N}\) 生成向量空间 \(V\)。
\item
  \begin{enumerate}
  \def\labelenumii{(\alph{enumii})}
  \tightlist
  \item
    设 G 是只有有限个连通分支的有限维交换实李群。G 紧的充要条件是：G 在复向量空间上的每个有限维解析线性表示都是半单的。（使用命题 32 的证明。）
  \end{enumerate}
\end{enumerate}

\begin{enumerate}
\def\labelenumi{(\alph{enumi})}
\setcounter{enumi}{1}
\tightlist
\item
  设 G 是只有有限个连通分支的有限维交换复李群。设 N 是 \(\exp_{G}\) 的核。
\end{enumerate}

N 在 C 上生成向量空间 \(L(G)\) 的充要条件是：G 的每个有限维解析线性表示都是半单的。（使用 (a)、引理 1 和命题 32。）

\begin{enumerate}
\def\labelenumi{\arabic{enumi}.}
\setcounter{enumi}{23}
\item
  李群 \(\mathbf{SL}(2,\mathbf{R})\) 是连通的且几乎单纯，但 \(\{I,-I\}\) 是 \(\mathbf{SL}(2,\mathbf{R})\) 的正规交换子群。
\item
  设 G 是几乎单纯的连通实或复李群，A 是 G 的正规子群。若 \(A \neq G\)，则 A 离散且为中心子群。（使用 §4 的练习 8。）因此，G 对其中心的商作为抽象群是单纯的。
\item
 设 G 是有限维连通实李群。假设 G 具有有限维单射连续线性表示 \(\rho\)。则 (G, G) 在 G 中闭。（设 R 是 G 的根基，S 是 G 的一个极大半单积分子群。利用练习 7 (b)，将问题化归为 G 是 S 与 R 的半直积的情形。根据第 I 章 § 6 命题 6，\(\rho\) 在 (G, R) 上是幺幂的。因此 \(\rho((G, R))\) 在线性群中闭，从而 (G, R) 在 G 中闭。）
\item
  设 G 是有限维可解单连通实李群，N 是 G 的最大连通幂零正规子群。则 G 具有一个在 N 上为幺幂的有限维单射连续线性表示。（对 G 的维数归纳。使用命题 20 以及第 I 章 § 7 定理 1。）†
\end{enumerate}

¶ 28. (a) 设 k 是特征为 0 的交换域，L 是 k 上的 n 维李代数，\(L_{1}\) 是一个 \((n-1)\) 维子代数，不含 L 的非零理想，且 \(a_{0}\) 是 L 中不属于 \(L_{1}\) 的元素。对于 \(i=2,3,\ldots\)，递归定义 \(L_{i}\) 为所有 \(x\in L_{i-1}\) 且满足 \([x,a_{0}]\in L_{i-1}\) 的元素组成的集合；记 \(L_{i}=L\) 对于 \(i\leqslant0\)。证明 \([L_{0},L_{i}]\subset L_{i+j}\)（对 \(i+j\) 归纳），然后证明对于 \(0\leqslant i\leqslant n\)，\(L_{i}\) 是 L 的余维为 i 的子代数（对 i 归纳）。

\begin{enumerate}
\def\labelenumi{(\alph{enumi})}
\setcounter{enumi}{1}
\tightlist
\item
  对于 0 \textless{} i \textless{} n，在 \(L_{i}\) 中选取元素 \(a_{i}\)，使其不属于 \(L_{i+1}\) 且满足 \([a_{0}, a_{i}] \equiv ia_{i-1} \pmod{L_{i}}\)。证明，对按字典序排列的有序对 \((i, j)\) 归纳，对于 \(0 \leqslant i < j < n\)、\(i + j - 1 < n\)，有
\end{enumerate}

\[
[ a _ {i}, a _ {j} ] \equiv (j - i) a _ {i + j - 1} (\mathrm{mod}. \mathrm{L} _ {i + j}).
\]

\begin{enumerate}
\def\labelenumi{(\alph{enumi})}
\setcounter{enumi}{2}
\item
  推出 L 或为 1 维，或为 2 维且非交换，或同构于 \(\mathfrak{sl}(2,k)\)。
\item
  设 A 是有限维实李代数。若 A 的子代数 B 存在满足 \(B_{1} \supset B\) 且 dim \(B_{1} = \dim B + 1\) 的子代数，则称 B 可扩张。记 \(R'(A)\) 为 A 的所有不可扩张子代数的交。记 \(R(A)\) 为 A 的最大理想，使其作为 A-模具有一个所有商模维数均为 1 的合成列。
\end{enumerate}

证明 \(\mathbf{R}'(\mathbf{A})\) 是 \(\mathbf{A}\) 的特征理想。（注意 \(\mathbf{R}'(\mathbf{A})\) 在 \(\operatorname{Aut}(\mathbf{A})\) 下稳定。）

\begin{enumerate}
\def\labelenumi{(\alph{enumi})}
\setcounter{enumi}{4}
\item
  证明 \(\mathbf{R}'(\mathbf{A})\supset \mathbf{R}(\mathbf{A})\)，且 \(\mathrm{R}'(\mathrm{A} / \mathrm{R}(\mathrm{A})) = \mathrm{R}'(\mathrm{A}) / \mathrm{R}(\mathrm{A})\)。
\item
  证明若 B 是 A 的子代数，则 R'(B) ⊃ R'(A) ∩ B。
\item
  若 A 可解，则 \(\mathrm{R}^{\prime}(\mathrm{A}) = \mathrm{R}(\mathrm{A})\)。（根据 (e)，可假设 \(\mathrm{R}(\mathrm{A}) = \{0\}\)。设 \(\mathrm{R}^{\prime}(\mathrm{A}) \neq \{0\}\)。由 (d)，存在包含于 \(\mathrm{R}^{\prime}(\mathrm{A})\) 中的 A 的极小非零理想 I。利用 (f) 及第 I 章 §5 定理 1 的推论 1，证明 dim I = 1，从而得到矛盾。）
\item
  设 P 是 \(\mathrm{R}^{\prime}(\mathrm{A})\) 的根基，B 是 A 的半单子代数，且 \(C = P + B\) 根据 (d) 是 A 的子代数。利用 (f)，证明相对于 P 的适当基，\(ad_{P}B\) 可表示为三角形。推出 \([P, B] = \{0\}\)。
\end{enumerate}

\((k)\) R(A) 是 \(\mathrm{R}'(\mathrm{A})\) 的根基。（使用 \((e)\)、\((f)\)、\((g)\)、\((h)\) 以及 A 的 Levi 分解。）设 \(\mathrm{R}'(\mathrm{A}) = \mathrm{R}(\mathrm{A}) \oplus \mathrm{T}\) 是 \(\mathrm{R}'(\mathrm{A})\) 的 Levi 分解。根据 \((h)\)，\(\mathrm{R}'(\mathrm{A}) = \mathrm{R}(\mathrm{A}) \times \mathrm{T}\)。于是 \(\mathrm{R}(\mathrm{T}) = \{0\}\)。将 \((c)\) 应用于 T 的单因子，推出 \(\mathrm{T} = \{0\}\)。因此已经证明 \(\mathrm{R}(\mathrm{A}) = \mathrm{R}'(\mathrm{A})\)。

\begin{enumerate}
\def\labelenumi{(\alph{enumi})}
\setcounter{enumi}{11}
\item
  设 G 是李代数为 A 的连通实李群。令 \(\mathcal{S}(G)\) 为 G 的子群 N 的集合，其中存在具有下列性质的递减子群序列 \((\mathrm{N}_m, \mathrm{N}_{m-1}, \ldots, \mathrm{N}_0)\)：\(\mathrm{N}_m = \mathrm{N}\)，\(\mathrm{N}_0 = \{\varepsilon\}\)，每个 \(\mathrm{N}_i\) 都是 G 的正规连通李子群，且 \(\dim \mathrm{N}_i / \mathrm{N}_{i-1} = 1\) 对所有 \(i > 0\) 成立。证明，李代数为 R(A) 的 G 的积分子群 R(G) 是 \(\mathcal{S}(G)\) 的最大元。（对 \(\dim \mathrm{R(A)}\) 归纳。设 I 是 A 的 1 维理想，N 是 G 中相应的积分子群。对 N 取商。若 N 不闭，使用 § 6 的练习 21 (a)。）
\item
  若 G 的李子群 H 存在满足 \(H_{1} \supset H\) 且 \(\dim H_{1} = \dim H + 1\) 的李子群，则称 H 可扩张。记 \(R'(G)\) 为 G 的所有不可扩张连通李子群的交。
\end{enumerate}

设 B 是 A 的不可扩张子代数。证明 G 中相应的积分子群是闭的。（使用命题 5。）推出

\[
\mathrm{L} (\mathrm{R} ^ {\prime} (\mathrm{G})) \subset \mathrm{R} (\mathrm{A}).
\]

从而 \(\mathrm{R}^{\prime}(\mathrm{G})\subset\mathrm{R}(\mathrm{G})\)。

\begin{enumerate}
\def\labelenumi{(\alph{enumi})}
\setcounter{enumi}{13}
\tightlist
\item
  证明 \(\mathrm{R}(\mathrm{G}) = \mathrm{R}'(\mathrm{G})\)。（将其化归为 \(\mathrm{R}'(\mathrm{G}) = \{e\}\) 的情形。然后使用 §6 的练习 22。）†
\end{enumerate}

\begin{enumerate}
\def\labelenumi{\arabic{enumi}.}
\setcounter{enumi}{28}
\tightlist
\item
  实李群 G 称为 (N) 型，如果它是有限维的，
\end{enumerate}

幂零、连通且单连通。若 G 是这样的群且 g 是其李代数，则当赋予 g 以由豪斯多夫法则定义的群结构时，映射 exp: g → G 是同构（参见第 II 章 §6 第 5 号注记 3）。记 log: G → g 为逆同构。

\begin{enumerate}
\def\labelenumi{(\alph{enumi})}
\item
  设 V 是 g 的一个向量 Q-子空间。证明下列条件等价：
\item
  V 是 \(\mathbf{Q}\) 上的 \(\mathfrak{g}\) 的李子代数；
\end{enumerate}

\begin{enumerate}
\def\labelenumi{(\roman{enumi})}
\setcounter{enumi}{1}
\tightlist
\item
  \(\exp(V)\) 是 G 的子群。
\end{enumerate}

（使用第 II 章 § 6 的练习 5。）

G 的子群 H 要具有 \(\exp(V)\) 的形式（其中 V 是 g 的向量 Q-子空间），充要条件是 H 在 G 中饱和（第 II 章 §4，练习 14），即关系 \(x \in G\)、\(x^{n} \in H\)、\(n \neq 0\) 蕴含 \(x \in H\)（同上）。若此条件满足，证明 H 是 G 的积分子群，当且仅当 \(\log(H)\) 是 g 的李 R-子代数。

\begin{enumerate}
\def\labelenumi{(\alph{enumi})}
\setcounter{enumi}{1}
\item
  设 V 是 g 的一个有限维李 Q-子代数，\((e_{1},\ldots,e_{m})\) 是 V 在 Q 上的一个基，\(\Lambda\) 是由 \(e_{1},\ldots,e_{m}\) 生成的 V 的子群。利用豪斯多夫法则是多项式这一事实，证明存在一个整数 \(d\geqslant1\)，使得对于 d 的每个非零倍数整数 r，\(\exp(r\Lambda)\) 都是 G 的子群。设 d 为这样的整数，r 为 d 的倍数。证明 \(\exp(r\Lambda)\) 离散，当且仅当 \((e_{1},\ldots,e_{m})\) 是 R 上的自由族，即 V \(\otimes_{q}R\) 到 g 的典范映射是单射。假设此条件成立；证明 G/exp(r \(\Lambda\) ) 紧，当且仅当 V \(\otimes_{q}R\to g\) 是双射，即 V 是 g 的一个 Q-型。（为证明充分性，对 g 的幂零类归纳。）
\item
  反之，设 \(\Gamma\) 是 G 的离散子群，令 \(\Gamma\) 为其在 G 中的饱和（第 II 章，同上），并令 \(\mathfrak{g}_{\Gamma} = \log (\Gamma) = \mathbf{Q}. \log (\Gamma)\) 为相应的李 \(\mathbf{Q}\)-子代数。假设 \(\mathbf{R}, \mathfrak{g} = \mathfrak{g}\)，即 \(\Gamma\) 不包含于 G 的任何不同的积分子群中。设 \(z\) 是满足 \(\neq 1\) 的 \(\Gamma\) 的中心元素，令 \(x = \log (z)\)。证明 \(x\) 属于 \(\mathfrak{g}\) 的中心。若 \(\mathrm{X} = \exp (\mathbf{R}x)\)，证明 \(\mathrm{X} / (\Gamma \cap \mathrm{X})\) 是紧的，且 \(\Gamma\) 在 \(\mathrm{G} / \mathrm{X}\) 中的像是离散的。通过对 dim G 归纳，推出 \(\mathrm{G} / \Gamma\) 紧，且 \(\mathfrak{g}_{\Gamma}\) 是 \(\mathfrak{g}\) 的 G-型。若 \(\Lambda\) 是 \(\mathfrak{g}_{\Gamma}\) 的格，证明存在整数 \(d \neq 0\)，使得 \(\Gamma\) 包含 \(\exp (d\Lambda)\)，并且 \(\exp (d\Lambda)\) 在 \(\Gamma\) 中的指数有限。
\end{enumerate}

设 H 是李代数为 \(\mathfrak{h}\) 的 G 的积分子群。证明 \(\mathrm{H} / (\mathrm{H} \cap \Gamma)\) 紧，当且仅当 \(\mathfrak{h}\) 关于 \(\mathbf{Q}\)-结构 \(g_{\Gamma}\) 是有理的（特别地，当 \(\mathfrak{h}\) 是 \(g\) 的下中心列或上中心列中的一项时成立）。

\begin{enumerate}
\def\labelenumi{(\alph{enumi})}
\setcounter{enumi}{3}
\item
  设 \(\Gamma\) 是 \(G\) 的离散子群，令 \(H\) 是 \(G\) 的最小积分子群，包含 \(\Gamma\)，且 \(\mathfrak{h}\) 是其李代数。证明 \(H / \Gamma\) 紧，且 \(\mathfrak{g} \otimes_{\mathbf{q}} \mathbf{R} \to \mathfrak{g}\) 是单射，其像为 \(\mathfrak{h}\)。（将 (c) 应用于幂零群 \(H\)。）
\item
  设 \(\Gamma\) 是 \(G\) 的离散子群。证明存在基 \((x_{1},\ldots ,x_{q})\) 作为 \(L(G)\) 的基，具有下列性质：
\item
  对所有 \(i \in [1, q]\)，\(\mathbf{R}x_i + \cdots + \mathbf{R}x_q\) 是记为 \(\mathfrak{n}_i\) 的 \(\mathbf{R}x_{i-1} + \cdots + \mathbf{R}x_q\) 的理想；
\end{enumerate}

\begin{enumerate}
\def\labelenumi{(\roman{enumi})}
\setcounter{enumi}{1}
\tightlist
\item
  存在 \(p \in (1, q]\)，使得 \(\Gamma\) 恰为如下乘积的集合
\end{enumerate}

\[
\exp (m _ {p} x _ {p}) \exp (m _ {p + 1} x _ {p + 1}) \dots \exp (m _ {q} x _ {q})
\]

其中 \(m_{p}, \ldots, m_{q}\) 属于 Z；

\begin{enumerate}
\def\labelenumi{(\roman{enumi})}
\setcounter{enumi}{2}
\tightlist
\item
  若 \(G / \Gamma\) 紧，则 \(\{n_1, n_2, \ldots, n_q\}\) 包含 \(g\) 的下中心列。（利用 (d) 和命题 16，将其化归为 \(G / \Gamma\) 紧的情形。然后利用 (c) 对 \(\dim G\) 归纳。）
\end{enumerate}

\begin{enumerate}
\def\labelenumi{(\arabic{enumi})}
\setcounter{enumi}{29}
\tightlist
\item
  举出一个 7 维实幂零李代数的例子，使其不存在一个基，关于该基的结构常数为有理数（参见第 I 章 § 4 的练习 18）。推出相应的 (N) 型群（练习 29）没有带紧子群的离散子群。（使用练习 29。）
\end{enumerate}

\begin{enumerate}
\def\labelenumi{\arabic{enumi}.}
\setcounter{enumi}{30}
\item
  设 G 和 \(G'\) 是两个 (N) 型李群（练习 29），\(\Gamma\) 是 G 的离散子群且 \(G/\Gamma\) 紧。证明每个同态 \(f: \Gamma \to G'\) 都能唯一延拓为从 G 到 \(G'\) 的李群态射（首先将 f 延拓到 \(\tilde{\Gamma}\) 的饱和 \(\Gamma\)，并导出从李 Q-代数 \(\log(\tilde{\Gamma})\) 到 \(G'\) 的李代数的一个同态；参见练习 29）。
\item
  设 G 是 (N) 型李群（练习 29），\(\Gamma\) 是 G 的离散子群。证明下列条件等价：
\end{enumerate}

\begin{enumerate}
\def\labelenumi{(\alph{enumi})}
\item
  G/Γ 是紧的；
\item
  \(\mathrm{G} / \Gamma\) 的测度有限（相对于一个非零的 G-不变正测度）；
\item
  每个包含 \(\Gamma\) 的 G 的积分子群都等于 G。（使用练习 29。）
\end{enumerate}

¶ 33. 设 \(\Gamma\) 是一个群。证明下列条件等价：(a) \(\Gamma\) 是幂零、无挠且有限生成的；

\begin{enumerate}
\def\labelenumi{(\alph{enumi})}
\setcounter{enumi}{1}
\item
  存在一个包含 \(\Gamma\) 作为离散子群的 (N) 型李群（练习 29）；
\item
  存在一个 (N) 型李群 G（练习 29），包含 \(\Gamma\) 作为离散子群，且 \(G / \Gamma\) 紧。
\end{enumerate}

（\(\langle b\rangle\) 与 \(\langle c\rangle\) 的等价性由练习 29 得出。蕴含 \((\varepsilon)\Rightarrow \langle a\rangle\) 通过对 \(\dim (\mathrm{G})\) 归纳并使用练习 29 \((c)\) 的方法得到证明。为证明 \((a)\Rightarrow (c)\)，先证明该李 \(\mathbf{Q}\)-代数（附属于 \(\Gamma\) 的饱和）是有限维的（参见第 II 章 § 6 的练习 4）；然后取该代数与 \(\mathbf{R}\) 的张量积以及相应的 (N) 型群。）

\begin{enumerate}
\def\labelenumi{\arabic{enumi}.}
\setcounter{enumi}{33}
\tightlist
\item
  设 \(G\) 是有限维连通实或复李群，\(\rho\) 是 \(G\) 在有限维复向量空间 \(V\) 上的解析线性表示。假设 \(\rho\) 是半单的。群 \(G\) 通过自同构作用于代数 \(S(V^*)\)，即 \(V\) 上的多项式函数代数。
\end{enumerate}

\begin{enumerate}
\def\labelenumi{(\alph{enumi})}
\item
  设 \(\mathsf{S}(\mathbf{V}^{*})^{\mathbb{G}}\) 是 \(\mathsf{S}(\mathbf{V}^{*})\) 的 G-不变元素的集合。存在投影算子 \(p\)，从 \(S(V^{*})\) 到 \(S(V^{*})^{\mathbb{G}}\)，它与 \(G\) 的运算交换，并使 \(G\)-稳定的 \(S(V^{*})\) 的每个向量子空间保持稳定。
\item
  设 I、J 是 \(S(V^*)\) 的在 G 下稳定的理想。令 A（分别为 B）是 I（分别为 J）在 V 中的零点集。假设 \(A \cap B = \varnothing\)。于是存在 \(u \in I\)，使 \(u\) 是 G-不变的，且在 B 上 \(u = 1\)。（根据《交换代数》第 7 章 §3 第 3 号命题 2，存在 \(v \in I\)、\(w \in J\)，使 \(v + w = 1\)。令 \(u = pv\)。证明 \(u\) 具有所需性质。）
\end{enumerate}

\begin{enumerate}
\def\labelenumi{\arabic{enumi}.}
\setcounter{enumi}{34}
\tightlist
\item
  设 \(\mathbf{G}\) 是 \(\mathbf{GL}(n,\mathbf{R})\) 的紧子群。
\end{enumerate}

\begin{enumerate}
\def\labelenumi{(\alph{enumi})}
\item
  设 A、B 是 \(\mathbf{R}^{\mathbf{n}}\) 中不交的紧 G-不变子集。存在 G-不变的多项式函数 \(u\)，定义在 \(\mathbf{R}^{\mathbf{n}}\) 上，使得 \(u = 1\) 在 \(\mathrm{A}\) 上成立且 \(u = 0\) 在 B 上成立。（存在连续实值函数 \(v\)，定义在 \(\mathbf{R}^{\mathbf{n}}\) 上，使得 \(v = 1\) 在 A 上成立且 \(v = 0\) 在 B 上成立。由 Stone-Weierstrass 定理，存在多项式函数 \(w\)，定义在 \(\mathbf{R}^{\mathbf{n}}\) 上，使得 \(|v - w| \leqslant \frac{1}{3}\) 在 \(\mathrm{A} \cup \mathrm{B}\) 上成立。通过相对于 G 的归一化 Haar 测度的积分过程得到 \(u\)。）
\item
  设 \(x \in \mathbf{R}^n\)。则 \(Gx\) 是 \(\mathbf{R}^n\) 中有限个 G-不变多项式的零点集。（使用 (a) 以及第 8 号定理 1 的推论 2。）
\end{enumerate}

\begin{enumerate}
\def\labelenumi{\arabic{enumi}.}
\setcounter{enumi}{35}
\tightlist
\item
  在 \(\mathbf{SL}(2, \mathbf{R})\) 中，能写成 \((a, b)\) 形式的元素（其中 \(a, b\) 属于 \(\mathbf{SL}(2, \mathbf{R})\)）正是满足 \(\neq -1\) 的元素。
\end{enumerate}

¶ 37. 设 G 是紧实李群，G' 是有限维连通实李群。假设 L(G) 是单纯的。设 \(\rho: G \to G'\) 是抽象群同态。假设 G 中存在 \(e_{\alpha}\) 的邻域 V，使得 \(\rho(V)\) 相对紧。则 \(\rho\) 连续。（设 V' 是 G' 中 \(e_{\alpha'}\) 的邻域。设 V'\,' \(\subset\) V' 是 \(e_{\alpha'}\) 的邻域，使得

\[
(x \in \mathrm{V} ^ {\prime \prime}, x _ {j} \in \rho (\mathrm{V}), y _ {j} \in \rho (\mathrm{V}) \text {   for   } 1 \leqslant j \leqslant n) \Rightarrow \left(\prod_ {i = 1} ^ {n} x _ {i} (y _ {i}, x) x _ {i} ^ {- 1} \in \mathrm{V} ^ {\prime}\right)
\]

其中 \(n = \dim G\)。可假设 \(\rho\) 非平凡。由练习 25，\(\operatorname{Ker} \rho\) 是有限的。因此 \(\rho(G)\) 不可数，因而不是离散的。于是存在 \(g \in G\)，其中心化子不开放，且 \(\rho(g) \in V''\)。使用 §4 练习 9 及其记号，有 \(\rho(M(g, n, V)) \subset V'\)，且 \(M(g, n, V)\) 是 G 中 \(e_G\) 的邻域。）

\begin{enumerate}
\def\labelenumi{\arabic{enumi}.}
\setcounter{enumi}{37}
\tightlist
\item
  设 \(G\) 是有限维实李群，\(G_0\) 是其单位元连通分支。考虑以下性质：
\end{enumerate}

\begin{enumerate}
\def\labelenumi{(\Alph{enumi})}
\setcounter{enumi}{5}
\tightlist
\item
  Ad(G) 在 Aut L(G) 中闭。
\end{enumerate}

证明在下列每种情形中 G 都具有性质 (F)：

\begin{enumerate}
\def\labelenumi{(\roman{enumi})}
\item
  G 连通且幂零；
\item
  L(G) 的每个导子都是内导子；
\item
  \(\mathrm{G}_0\) 具有性质 (F)，且 \(\mathrm{G} / \mathrm{G}_0\) 是有限的；
\item
  G 是上三角群；
\item
  \(G_{0}\) 是半单的。
\end{enumerate}

\begin{enumerate}
\def\labelenumi{\arabic{enumi}.}
\setcounter{enumi}{38}
\tightlist
\item
  设 G 是有限维连通实李群，H 是 G 的积分子群，\(\bar{H}\) 是 H 在 G 中的闭包。假设 H 具有性质 (F)（练习 38）。
\end{enumerate}

\begin{enumerate}
\def\labelenumi{(\alph{enumi})}
\item
  设 \(x \in \overline{\mathbf{H}}\)。则 \(\mathrm{Ad}_{\mathrm{L(G)}}x\) 保持 \(\mathrm{L(H)}\) 稳定（第 2 号，命题 5）；令 \(u(x)\) 为其在 \(\mathrm{L(H)}\) 上的限制。于是 \(u(\overline{\mathrm{H}}) = \mathrm{Ad}_{\mathrm{L(H)}}(\mathrm{H})\)。（注意 \(\mathrm{Ad}_{\mathrm{L(H)}}(\mathrm{H})\) 在 \(u(\overline{\mathrm{H}})\) 中稠密，并应用性质 (F)。）
\item
  设 C 是 \(\overline{\mathrm{H}}\) 的中心。则 \(\overline{\mathrm{H}} = \mathrm{C.H}\)，且 C 是 H 的中心的闭包。（根据 (a)，若 \(x \in \overline{\mathrm{H}}\)，则存在 \(y \in \mathrm{H}\)，使 \(\mathrm{Ad}_{\mathrm{L(H)}} x = \mathrm{Ad}_{\mathrm{L(H)}} y\)。于是 \(\mathrm{Ad}(x^{-1}y)\) 在 L(H) 上是恒等的，从而 \(x^{-1}y \in Z_{\overline{\mathrm{H}}}(\mathrm{H})\)，且 \(x \in Z_{\overline{\mathrm{H}}}(\mathrm{H})\)。H。由连续性，\(Z_{\overline{\mathrm{H}}}(\mathrm{H})\) 等于 C。群 C 在 \(\overline{\mathrm{H}}\) 中闭，因而在 G 中也闭；所以 \(\overline{\mathrm{C}} \cap \overline{\mathrm{H}} \subset \mathrm{C}\)。设 \(x \in \mathrm{C}\)。存在 H 中的序列 \((x_n)\)，使 \(x_n\) 趋于 \(x\)。于是 \(\mathrm{Ad}_{\mathrm{L(H)}} x_n\) 趋于 1，从而存在 H 中的序列 \((y_n)\)，使 \(y_n\) 趋于 \(e\)，且 \(\mathrm{Ad}_{\mathrm{L(H)}} y_n = \mathrm{Ad}_{\mathrm{L(H)}} x_n\)。于是 \(y_n^{-1}x_n \in \mathrm{C} \cap \mathrm{H}\)，且 \(y_n^{-1}x_n\) 趋于 \(x\)。）
\item
  \(\mathbf{H} = \overline{\mathbf{H}}\) 当且仅当 \(\mathbf{H}\) 的中心在 \(\mathbf{G}\) 中闭。（使用 (b)。）
\end{enumerate}

\begin{enumerate}
\def\labelenumi{\arabic{enumi}.}
\setcounter{enumi}{39}
\tightlist
\item
  设 \(\mathbf{H}\) 是有限维实李群，\(\mathbf{H}_0\) 是其单位元连通分支。
\end{enumerate}

\begin{enumerate}
\def\labelenumi{(\alph{enumi})}
\item
  假设 \(H_{0}\) 满足练习 38 的条件 (F)，\(H/H_{0}\) 有限，且 \(H_{0}\) 的中心紧。设 G 是有限维实李群，\(f:H\to G\) 是具有离散核的连续同态。则 \(f(\mathrm{H})\) 在 G 中闭。（将其化归为 H 连通的情形。由于 f 的核是离散的，\(f(\mathrm{H})\) 的中心是 H 的中心在 f 下的像（引理 1），因而在 G 中闭。应用练习 39 (c)。）
\item
  假设 \(H_{0}\) 是半单的且 \(H/H_{0}\) 有限。则 H 在有限维连续线性表示下的像在线性群中闭。（应用 (a) 和练习 7 (b)。）
\end{enumerate}

\begin{enumerate}
\def\labelenumi{\arabic{enumi}.}
\setcounter{enumi}{40}
\tightlist
\item
  设 \(G\) 是有限维连通实李群，\(\rho: G \to \mathbf{GL}(n, \mathbf{C})\) 是具有有限核的连续同态。则 \(\rho((G, G))\) 在 \(\mathbf{GL}(n, \mathbf{C})\) 中闭。（利用练习 7 (b) 和 40 (a)，将其化归为 \(G\) 是半单群 \(S\) 与其根基 \(R\) 的半直积的情形。\(\rho((G, R))\) 的每个元素都是幺幂的，故 \(\rho((G, R))\) 闭。向量子空间 \(V_1\) 由 \(\rho((G, R))\) 的不动点组成，满足 \(\neq \{0\}\)。它在 \(\rho(G)\) 下稳定。考虑 \(\mathbf{C}^n / V_1\)，对 \(n\) 归纳并使用 \(\rho(S)\) 的完全可约性，可以看出能够写成
\end{enumerate}

\[
\mathbf {C} ^ {n} = \mathrm{V} _ {1} \oplus \dots \oplus \mathrm{V} _ {q}
\]

其中每个 \(V_{i}\) 在 \(\rho(S)\) 下稳定，且对所有 i 有 \(\rho((G,R))(V_{i})\subset V_{1}\oplus\cdots\oplus V_{i-1}\)。另一方面，\(\rho(S)\) 是闭的（练习 40 (b)）。最后，\(\rho((G,G))=\rho(S)\cdot\rho((G,R)).\)

\begin{enumerate}
\def\labelenumi{\arabic{enumi}.}
\setcounter{enumi}{41}
\tightlist
\item
  设 \(G\) 是有限维连通实李群。假设
\end{enumerate}

G 具有一个单射连续线性表示 \(\rho: \mathrm{G} \to \mathbf{GL}(n, \mathbf{C})\)。则 G 具有一个线性表示 \(\sigma: \mathrm{G} \to \mathbf{GL}(n', \mathbf{C})\)，它把 G 同态到 \(\mathbf{GL}(n', \mathbf{C})\) 的一个闭子群上。（根据练习 41，\(\rho((\mathrm{G}, \mathrm{G}))\) 在 \(\mathbf{GL}(n, \mathbf{C})\) 中闭，因而 (G, G) 在 G 中闭。令 \(\rho: \mathrm{G} \to (\mathrm{G}, \mathrm{G})\) 为典范态射，\(\tau\) 为一个像闭的单射线性表示，其像 \(\mathrm{G}/(\mathrm{G}, \mathrm{G})\) 连通且交换。取 \(\sigma = \rho \oplus (\tau \circ p)\)。）

\exercisesfor{10}

\begin{enumerate}
\def\labelenumi{\arabic{enumi}.}
\item
  设 \(G\) 为有限维连通实李群。证明，从 \(\operatorname{Aut}(G)\) 到 \(\operatorname{Aut}(L(G))\) 的典范映射一般不是满射（取 \(G = T\)）。
\item
  假设 \(\mathbf{K} = \mathbf{Q}_p\) 。设 \(\mathbf{G}\) 为有限维李群。证明下列条件等价：
\end{enumerate}

\begin{enumerate}
\def\labelenumi{(\alph{enumi})}
\item
  存在 \(x_{1}, x_{2}, \ldots, x_{n}\) 属于 \(G\)，使 \(G\) 的由 \(\{x_{1}, \ldots, x_{n}\}\) 生成的子群在 \(G\) 中稠密；
\item
  G 由一个紧子集生成。
\end{enumerate}

（要证明 (b) 蕴含 (a)，注意：若 \((e_1, \ldots, e_n)\) 是 \(\mathrm{L}(\mathrm{G})\) 的一个基，则 \((\exp \mathbf{Z}_p e_1)(\exp \mathbf{Z}_p e_2) \ldots (\exp \mathbf{Z}_p e_n)\) 是 \(e\) 的一个邻域。）

\begin{enumerate}
\def\labelenumi{\arabic{enumi}.}
\setcounter{enumi}{2}
\tightlist
\item
  设 \(G\) 为满足 \((x, y) \in \mathbf{Q}_p \times \mathbf{Q}_p\) 且 \(|y| \leqslant 1\) 的点的集合。它是 \(\mathbf{Q}_p \times \mathbf{Q}_p\) 的一个开子群。然后 \(L(G) = \mathbf{Q}_p \times \mathbf{Q}_p\)。设 \(\alpha\) 为无穷小自同构 \((x, y) \mapsto (0, x)\)。于是命题 3 的结论是错误的。
\end{enumerate}

\begin{enumerate}
\def\labelenumi{(\arabic{enumi})}
\setcounter{enumi}{3}
\tightlist
\item
  设 \(\mathbf{K}\) 是 \(\mathbf{Q}_p\) 的二次扩张，且 \(\omega \in \mathbf{K} - \mathbf{Q}_p\)。令 \(\mathrm{G} = \mathbf{Q}_p + \mathbf{Z}_p\omega\)，将其视为 \(\mathbf{K}\) 上的李群。 \(\mathrm{G}\) 的自同构全是 \((x,y)\mapsto (\lambda x,y)\)，其中 \(\lambda\) 在 \(\mathbf{Z}_p\) 中可逆。这些自同构不能在 \(\mathbf{K}\) 上被赋予具有定理 1 所述性质的李群结构。
\end{enumerate}

\specialchapter{历史注记（第 I 至 III 章）}

\hnsection{起源}

近一个世纪以来被称为“李群理论”的理论，基本上由一位数学家 Sophus Lie 发展起来。

在开始叙述这段历史之前，我们先简要概述为此铺平道路的若干早期研究。

\textbf{(a) 变换群（Klein-Lie，1869--1872）}

约在 1860 年，一个有限集合的置换群理论正在发展，并开始得到应用（Serret、Kronecker、Mathieu、Jordan）。另一方面，正蓬勃发展的不变量理论使数学家熟悉了某些在复合下稳定的几何变换的无限集合（尤其是线性变换或射影变换）。然而，在 Jordan {[}7{]} 于 1868 年研究“运动群”（三维欧氏空间位移群的闭子群）之前，似乎还没有人有意识地把这两股思想联系起来。

1869 年，年轻的 Felix Klein（1849--1925）在柏林与挪威人 Sophus Lie（1842--1899）结为朋友；Lie 比他年长几岁，二人因共同关心 Plücker 的“直线几何”，尤其是直线复形理论，而走到了一起。大约就在这时，Lie 形成了他最具原创性的思想之一：把不变量的概念引入分析学和微分几何；其思想来源之一，是他注意到微分方程的经典“求积”积分方法完全依赖于方程在一个“连续”变换族下保持不变。1869 年的第一篇著作（由 Klein 编订）中，Lie 运用了这一思想；他在那里研究了“Reye 复形”（由四条直线组成的集合：这些直线在四面体的各面上截出四个具有给定交比的点）以及以该复形的直线为切线的曲线和曲面 {[}3a{]}：他的方法依赖于 Reye 复形在保持四面体顶点不变的三参数交换群（PGL(4, C) 的极大环面）下的不变性。Klein 和 Lie 于 1870 年春在巴黎共同写作的论文中也贯穿着同一思想 {[}1, a{]}；在那里，他们基本确定了平面射影群 PGL(3, C) 的交换连通子群，并研究了其轨道的几何性质（称为曲线或曲面 V）；通过统一的步骤，这使他们获得了各种代数或超越曲线的性质，例如 \(y = cx^m\) 或对数螺线。两篇论文都突出了 Galois 和 Jordan 的理论给他们留下的深刻印象（Jordan 关于 Galois 的评论于 1869 年发表在 Math. Annalen；另一方面，Lie 早在 1863 年就听说过 Galois 理论）。Klein 于 1871 年开始对非欧几何感兴趣，并由此开始研究一种对所有已知几何进行分类的原则，最终于 1872 年形成“Erlangen 纲领”。至于 Lie，他在 1873 年致 A. Mayer 的一封信中（{[}3{]}，第 V 卷，第 584 页）把自己关于变换群的思想追溯到巴黎逗留时期；在 1871 年的一篇著作中（{[}3 b{]}，第 208 页），他已经使用“变换群”一词，并明确提出确定 GL(n, C) 的所有子群（“连续或不连续”）的问题。坦率地说，Klein 和 Lie 在进入这个新的数学世界时都遇到了一些困难，Klein 把 Jordan 新近出版的“论著”说成是一本“七印封缄的书”（{[}2{]}，第 51 页）；此外，他在谈到 {[}1 a{]} 和 {[}1 b{]} 时写道：“关于连续算子群这一启发性思想的一切功劳都应归于 Lie，特别是微分方程和偏导数的积分。后来他在连续群理论中发展出的所有概念，虽然都已以萌芽形式存在，但阐述得如此不充分，以至于只有经过长时间交谈，我才能使他相信许多细节，例如首先要相信曲线 V 的确存在”（{[}2{]}，第 415 页）。

\textbf{(b) 无穷小变换}

“无穷小”变换的概念至少可以追溯到无穷小分析学的开端；我们知道，笛卡尔承认“在无穷小范围内”每个平面运动都可视为旋转时，发现了瞬时旋转中心：18 世纪分析力学的发展完全建立在类似思想上。1851 年，Sylvester 试图构造线性群 \(\mathbf{GL}(3,\mathbf{C})\) 及其某些子群的不变量，把其矩阵中出现的参数视为形如 \(\alpha_{j}dt\) 的“无穷小”增量；他把函数 \(f((z_{j}))\) 是不变量这一事实写成方程 \(f((z_{j}+\alpha_{j}(dt))=f((z_{j}))\)；这给出了关于 f 的线性偏微分方程 Xf=0，其中

\[
\mathrm{X} f = \sum_ {j} \alpha_ {j} \frac {\partial f}{\partial z _ {j}},\tag{1}
\]

其中 X 因而是一个微分算子，即“沿方向参数 \(\alpha_{j}\) 的导数”（{[}5{]}，第 3 卷，第 326、327 页）；Sylvester 似乎认为这是一个极其重要的一般原理，但此后似乎从未再回到这个问题。不久之后，Cayley（{[}6{]}，第 II 卷，第 164--178 页）以类似方式研究某些表示下 \(\mathbf{SL}(2,\mathbf{C})\) 的不变量，并证明它们是两个一阶偏微分方程 \(Xf=0\)、\(Yf=0\) 的解，其中 X 和 Y 如上由“无穷小”变换得到

\[
\left( \begin{array}{c c} 0 & 0 \\ d t & 0 \end{array} \right) \quad \text { and } \quad \left( \begin{array}{c c} 0 & d t \\ 0 & 0 \end{array} \right).
\]

用现代术语说，这表述为 X 和 Y 生成李代数 \(\mathfrak{sl}(2,\mathbf{C})\)；此外，Cayley 明确计算了括号 XY -- YX，并证明它也来自一个“无穷小”变换。

在他 1868 年关于运动群的论文 {[}7{]} 中，Jordan 自始至终使用“无穷小变换”这一概念，但完全是从几何观点出发的。他无疑提出了由无穷小变换“生成”的单参数群这一思想：对 Jordan 而言，它是通过“适当地重复”无穷小变换而得到的变换集合（同上，第 243 页）。Klein 和 Lie 在他们的论文中使用了同一表述“重复的无穷小变换” {[}1 b{]}，但上下文表明，他们所指的是一个微分系统的积分。如果他们考虑的单参数群由变换 \(x' = f(x, y, t)\)、\(y' = g(x, y, t)\) 组成，则相应的“无穷小变换”由下式给出：

\[
d x = p (x, y) d t, \quad d y = q (x, y) d t
\]

其中 \(p(x,y) = \frac{\partial f}{\partial t} (x,y,t_0), q(x,y) = \frac{\partial g}{\partial t} (x,y,t_0)\)，且 \(t_0\) 对应于

该群的恒等变换。由于 Klein 和 Lie 明确知道函数 f 和 g，他们很容易验证下列函数

\[
t \mapsto f (x, y, t) \quad \text { and } \quad t \mapsto g (x, y, t)
\]

用参数形式给出微分方程的积分曲线

\[
q (\xi , \eta) d \xi = p (\xi , \eta) d \eta
\]

经过点 \((x,y)\)，但他们没有给出一般论证；而且在其论文余下部分也从未使用这一事实。

\textbf{(c) 接触变换}

接下来的两年中，Lie 似乎放弃了变换群理论（尽管他仍与于 1872 年发表“纲领”的 Klein 保持非常密切的联系），转而研究接触变换、一阶偏微分方程的积分以及这两个理论之间的关系。这里不讨论这些问题的历史，只提及几点似乎在变换群理论的起源中发挥了重要作用的事实。

接触变换的概念同时推广了点变换和逆极变换。粗略地说，\(\mathbf{C}^{\mathrm{n}}\) 上的接触变换† 是一个开子集 \(\Omega\) 的同构，其中该开子集属于流形 \(T^{\prime}(\mathbf{C}^{\mathrm{n}})\)，该流形由指向 \(\mathbf{C}^{\mathrm{n}}\) 的余切向量组成；另一个开子集 \(\Omega^{\prime}\) 属于 \(T^{\prime}(\mathbf{C}^{\mathrm{n}})\)，且该同构把 \(\Omega\) 的典范 1-形式映到 \(\Omega^{\prime}\) 的典范 1-形式。换言之，若 \((x_{1},\ldots ,x_{n},p_{1},\ldots ,p_{n})\) 是 \(T^{\prime}(\mathbf{C}^{\mathrm{n}})\) 的典范坐标，则接触变换是同构 \((x_{1},p_{1})\mapsto (\mathbf{X}_{\mathrm{i}},\mathbf{P}_{\mathrm{i}})\)，满足关系 \(\sum_{i = 1}^{n}\mathbf{P}_{i}d\mathbf{X}_{i} = \sum_{i = 1}^{n}p_{i}dx_{i}\)。这类变换出现在研究形如

\[
\mathrm{F} \left(x _ {1}, x _ {2}, \dots , x _ {n}, \frac {\partial z}{\partial x _ {1}}, \dots , \frac {\partial z}{\partial x _ {n}}\right) = 0.\tag{2}
\]

Lie 在研究这些问题的过程中熟悉了 Poisson 括号的运算

\[
(f, g) = \sum_ {i = 1} ^ {n} \left(\frac {\partial f}{\partial x _ {i}} \frac {\partial g}{\partial p _ {i}} - \frac {\partial g}{\partial x _ {i}} \frac {\partial f}{\partial p _ {i}}\right)\tag{3}
\]

以及微分算子 \(\ddagger_{\xi}\) 的括号 {[}X, Y{]} = XY - YX；他把 Poisson 括号 (3) 解释为与 g 对应的 (1) 型变换对 f 的作用，并注意到，Poisson 括号的 Jacobi 恒等式意味着，与 f 和 g 对应的微分算子的括号与括号 (g, h) 相联系。对满足 (F, g) = 0 的函数 g 的研究（它出现在 Jacobi 积分偏微分方程 (2) 的方法中）使 Lie 转而研究保持给定方程不变的无穷小接触变换。最后，Lie 研究了这样的函数集 \((u_{j})_{1 \in j \leq m}\)：它们是 \(x_{i}\) 和 \(p_{i}\) 的函数，并且括号 \((u_{j}, u_{k})\) 是 \(u_{h}\) 的函数；他称这些集合为“群”（它们实质上已由 Jacobi 研究过） {[}3 c{]}。

\hnsection{连续群与无穷小变换}

1873 年秋，Lie 突然重新开始研究变换群，并取得决定性成果。根据他 1873--1874 年写给 A. Mayer 的几封信（{[}3{]}，第 5 卷，第 584--608 页），在能够追踪其思路的范围内，他从作用于 n 个变量的“连续群”变换出发

\[
x _ {i} ^ {\prime} = f _ {i} (x _ {1}, \dots , x _ {n}, a _ {1}, \dots , a _ {r}) \quad (1 \leqslant i \leqslant n)\tag{4}
\]

它有效地†依赖于 r 个参数 \(a_{1},\ldots,a_{r}\)；他注意到，若变换 (4) 在参数取 \(a_{1}^{0},\ldots,a_{r}^{0}\) 时是恒等变换‡，则 \(x_{i}\) 的一阶 Taylor 展开：

\[
\begin{array}{r l} f _ {i} (x _ {1}, \dots , x _ {n}, a _ {1} ^ {0} + z _ {1}, \dots , a _ {r} ^ {0} + z _ {r}) & \\ = x _ {i} + \sum_ {k = 1} ^ {r} z _ {k} X _ {k i} (x _ {1}, \dots , x _ {n}) + \dots & (1 \leqslant i \leqslant n) \end{array} \tag {5}
\]

给出一个关于 r 个参数 \(z_{j}\) 线性依赖的“通用”无穷小变换

\[
d x _ {i} = \left(\sum_ {k = 1} ^ {r} z _ {k} \mathrm{X} _ {k i} (x _ {1}, \dots , x _ {n})\right) d t \quad (1 \leqslant i \leqslant n).\tag{6}
\]

仿照他与 Klein 合写的论文，Lie 对微分系统作了积分

\[
\frac {d \xi_ {1}}{\sum_ {k} z _ {k} X _ {k 1} (\xi_ {1} , \dots , \xi_ {n})} = \dots = \frac {d \xi_ {n}}{\sum_ {k} z _ {k} X _ {k n} (\xi_ {1} , \dots , \xi_ {n})} = d t,\tag{7}
\]

这使他对每个点 \((z_{1},\ldots,z_{r})\) 得到一个单参数群

\[
t \mapsto x _ {i} ^ {\prime} = g _ {i} (x _ {1}, \dots , x _ {n}, z _ {1}, \dots , z _ {r}, t) \quad (1 \leqslant i \leqslant n)\tag{8}
\]

从而 \(g_{i}(x_{1},\ldots,x_{n},z_{1},\ldots,z_{r},0)=x_{i}\) 对所有 i 成立。他用一种巧妙的方法，利用变换 (4) 在复合下构成稳定集合这一事实，证明单参数群 (8) 是给定群的子群 {[}3 d{]}。整个理论的关键新思想，是把函数 (4) 的 Taylor 展开取到二阶。他的论证过程相当混乱而带有启发式色彩（{[}3 d{]} 和 {[}3{]}，第 5 卷，第 600--601 页）；可以叙述如下。对于足够小的 \(z_{r}\)，令 (8) 中 t=1；于是得到该群变换的新参数 \(z_{1},\ldots,z_{r}\)（这实际上是“典范参数”一词首次出现）。然后根据 (7)，

\[
\frac {\partial g _ {i}}{\partial t} = \sum_ {k} z _ {k} \mathrm{X} _ {k i} (x _ {1} ^ {\prime}, \dots , x _ {n} ^ {\prime})
\]

由此

\[
\begin{array}{r l} \frac {\partial^ {2} g _ {i}}{\partial t ^ {2}} & = \sum_ {k, j} z _ {k} \frac {\partial \mathrm{X} _ {k i}}{\partial x _ {j}} (x _ {1} ^ {\prime}, \ldots , x _ {n} ^ {\prime}) \frac {\partial x _ {j} ^ {\prime}}{\partial t} \\ & = \sum_ {k, j} z _ {k} \frac {\partial \mathrm{X} _ {k i}}{\partial x _ {j}} (x _ {1} ^ {\prime}, \ldots , x _ {n} ^ {\prime}) \left(\sum_ {n} z _ {h} \mathrm{X} _ {h j} (x _ {1} ^ {\prime}, \ldots , x _ {n} ^ {\prime})\right) \end{array}
\]

从而得到

\[
\begin{array}{r} x _ {i} ^ {\prime} = x _ {i} + \Big (\sum_ {k} z _ {k} \mathrm{X} _ {k i} (x _ {1}, \ldots , x _ {n}) \Big) t \\ + \frac {1}{2} \Big (\sum_ {k, n, j} z _ {k} z _ {h} \frac {\partial \mathrm{X} _ {k i}}{\partial x _ {j}} (x _ {1}, \ldots , x _ {n}) \mathrm{X} _ {h j} (x _ {1}, \ldots , x _ {n}) \Big) t ^ {2} + \dots , \end{array}
\]

由此，当 \(t = 1\) 时，关于参数 \(z_{j}\) 的 Taylor 展开为

\[
x _ {i} ^ {\prime} = x _ {i} + \left(\sum_ {j} z _ {k} \mathrm{X} _ {k i}\right) + \frac {1}{2} \left(\sum_ {k, h, j} z _ {k} z _ {h} \mathrm{X} _ {h j} \frac {\partial \mathrm{X} _ {k i}}{\partial x _ {j}}\right) + \dots (1 \leqslant i \leqslant n).\tag{9}
\]

将这些关系简写为 \(x' = \mathbf{G}(x, z)\)，用来联系下列向量

\[
x = (x _ {1}, \dots , x _ {n}), \quad x ^ {\prime} = (x _ {1} ^ {\prime}, \dots , x _ {n} ^ {\prime}), \quad z = (z _ {1}, \dots , z _ {r});
\]

这些变换集合关于复合的基本稳定性可写成

\[
\mathrm{G} (\mathrm{G} (x, u), v) = \mathrm{G} (x, \mathrm{H} (u, v))\tag{10}
\]

其中 \(\mathbf{H} = (\mathbf{H}_1, \ldots, \mathbf{H}_r)\) 与 \(x\) 无关；显然 \(\mathrm{H}(u, 0) = u\) 且 \(\mathrm{H}(0, v) = v\)，从而展开式为

\[
\mathrm{H} _ {i} (u, v) = u _ {i} + v _ {i} + \frac {1}{2} \sum_ {h, k} c _ {i k h} u _ {h} v _ {k} + \dots ,\tag{11}
\]

其中省略的项不关于 u 或 v 线性。将 (10) 与 (9)、(11) 结合变换，然后比较两边关于 \(u_{h}v_{k}\) 的项，Lie 得到关系

\[
\sum_ {j = 1} ^ {n} \left(\mathrm{X} _ {h j} \frac {\partial \mathrm{X} _ {k i}}{\partial x _ {j}} - \mathrm{X} _ {k j} \frac {\partial \mathrm{X} _ {h i}}{\partial x _ {j}}\right) = \sum_ {l = 1} ^ {r} c _ {l h k} \mathrm{X} _ {l i} (1 \leqslant h, k \leqslant r, 1 \leqslant i \leqslant n). \tag{12}
\]

他在偏微分方程理论方面的经验使他把这些条件写成更简单的形式：按照 (1) 的模式，对 (6) 中令 \(z_{k} = 1\)、\(z_{h} = 0\)（\(h \neq k\)）所得到的每个无穷小变换，赋予微分算子

\[
\mathrm{A} _ {k} (f) = \sum_ {i = 1} ^ {n} \mathrm{X} _ {k i} \frac {\partial f}{\partial x _ {i}},\tag{13}
\]

并将条件 (12) 改写为

\[
[ \mathrm{A} _ {h}, \mathrm{A} _ {k} ] = \sum_ {l} c _ {l h k} \mathrm{A} _ {l},\tag{14}
\]

这成为他理论的基石。在此之前，他不加区分地使用“无限小变换”和“无穷小变换”这两个术语（例如 {[}3 e{]}）；关系 (14) 的简洁性使他把无穷小变换 \(dx_{i} = X_{ki}dt\)（\(1 \leqslant i \leqslant n\)）的算子 (13) 称为“符号” {[}3 d{]}，很快又把算子 (13) 本身称为“无穷小变换”（{[}3 d{]} 和 {[}3{]}，第 5 卷，第 589 页）。

随后，他意识到“连续群”理论与自己早先关于接触变换和偏微分方程的研究之间存在密切联系。这种汇合使他充满热情：1874 年他在写给 Mayer 的信中说：“我早期的工作仿佛早已在那里，等待着创立变换群的新理论”（{[}3{]}，第 5 卷，第 586 页）。

在接下来的几年里，Lie 继续研究变换群。除下文（§ III）概述的一般定理外，他还得到了一些更特殊的结果：确定直线和平面的变换群、射影群中的低余维子群、至多含 6 个参数的群，等等。他并没有长期放弃微分方程。事实上，对他而言，变换群理论似乎是积分微分方程的一种工具，其中变换群所起的作用类似于代数方程的 Galois 群。†

我们还要指出，这项研究也促使他引入某些具有无穷多个参数的变换集合，他称之为“无限连续群” \(^{‡}\)；他把上面 (4) 型、具有有限个参数的变换群称为“有限连续群”。

\hnsection{李群—李代数“词典”}

Lie 自 1874 年起在大量论文中发展出的“有限连续群”理论，在他与 F. Engel§ 合作的宏大著作“变换群理论”（{[}4{]}，1888--1893{]}）中得到系统阐述；该理论是第一卷和第三卷最后五章的研究对象，第二卷则专门讨论接触变换。

正如标题所示，这部著作只关注式 (4) 意义下的变换群，其中“变量”空间 \(x_{i}\) 和“参数”空间 \(a_{i}\) 起初扮演着如此重要的角色。另一方面，当时“抽象”群的概念尚未被清楚地分离出来；Lie 在 1885 年（{[}3/j{]}，§5）注意到，在 (10) 的记号下，给出群中两个变换复合参数的方程 \(w = \mathrm{H}(u, v)\) 定义了一个新群，他把它看作参数空间上的变换群，从而得到他所谓的“参数群”（他甚至得到两个，这正是左平移群和右平移群†）。

式 (4) 中的变量 \(x_{i}\) 和参数 \(a_{i}\) 原则上假定为复数（第三卷第 XIX--XXIV 章除外），函数 \(f_{i}\) 假定为解析的；Lie 和 Engel 当然知道，这些函数一般并非对 \(x_{i}\) 和 \(a_{i}\) 的所有复值都有定义，因而这类函数的复合会产生严重困难（{[}4{]}，第 1 卷，第 15--17、33--40 页及各处）；尽管他们在全文中几乎总是写得仿佛所研究的变换可以无条件复合，这无疑是为了论述方便，但在必要时他们明确重申了“局部”观点（例如参见同上第 168 或 189 页，或同书第 3 卷第 2 页页下注）；换言之，他们研究的数学对象接近于本书所谓的运算律片段。他们有时也考虑全局群，例如经典群的 4 个系列（{[}4{]}，第 3 卷，第 682 页），但似乎没有问过一般而言什么构成“全局群”；他们满足于为经典群的“参数”（这些群的“变量”并不造成困难，因为所涉及的变换是 \(\mathbf{C}^n\) 的线性变换）取得恒等变换邻域中的“局部”参数系，而不去考虑所写公式的有效域。不过，他们提出了一个从局部理论自然产生的问题：研究“混合”群，即具有有限个连通分支的群，例如正交群（{[}4{]}，第 1 卷，第 7 页）。他们把这一研究表述为：研究一个在复合及取逆下稳定的变换集合，它是若干集合 \(\mathrm{H}_j\) 的并；每个集合由如 (4) 中的函数系统（\(f_i^{(D)}\)）描述；每个 \(\mathrm{H}_j\) 的（本质）参数个数起初也假定依赖于 \(j\)，但他们证明事实上这个数对所有 \(\mathrm{H}_j\) 都相同。他们的主要结果于是是：存在一个有限连续群 G，使 \(\mathrm{H}_j = \mathrm{G} \circ h_j\) 对某个 \(h_j \in \mathrm{H}_j\) 及所有 \(j\) 成立；他们还证明 G 在混合群中正规，并指出，后者不变量的确定归结为 G 的不变量和一个不连续群的不变量（{[}4{]}，第 1 卷，第 18 章）。

在 {[}4{]} 中发展起来的一般理论最终形成了一部“词典”（作者们没有非常系统地如此表述），把“有限连续群”的性质翻译为其无穷小变换集合的性质。它建立在“Lie 的三个定理”之上，每个定理都由一个断言及其逆命题组成。

第一定理（{[}4{]}，第 1 卷，第 33、72 页以及第 3 卷，第 563 页）首先断言：若 (4) 中的参数是本质的，则函数 \(f_{i}\) 满足如下形式的偏微分方程组

\[
\frac {\partial f _ {i}}{\partial a _ {j}} = \sum_ {k = 1} ^ {r} \xi_ {k i} (f (x, a)) \psi_ {k j} (a) \quad (1 \leqslant i \leqslant n)\tag{15}
\]

其中矩阵 \((\xi_{kl})\) 秩最大且 \(\det(\psi_{kj}) \neq 0\)；反之，若函数 \(f_{i}\) 具有这一性质，则公式 (4) 定义一个变换的群芽。

第二定理（{[}4{]}，第 1 卷，第 149、158 页以及第 3 卷，第 590 页）给出一方面的 \(\xi_{ki}\) 与另一方面的 \(\psi_{ij}\) 之间的关系：关于 \(\xi_{ki}\) 的条件可写成

\[
\sum_ {k = 1} ^ {n} \left(\xi_ {i k} \frac {\partial \xi_ {f l}}{\partial x _ {k}} - \xi_ {j k} \frac {\partial \xi_ {t l}}{\partial x _ {k}}\right) = \sum_ {k = 1} ^ {r} c _ {i j} ^ {k} \xi_ {k l} (1 \leqslant i, j \leqslant r, 1 \leqslant l \leqslant n)\tag{16}
\]

其中 \(c_{ii}^{k}\) 是常数 \((1\leqslant i,j,k\leqslant r)\)，关于 \(i,j\) 斜对称。Maurer 给出的 \(\psi_{ij}\) 条件（{[}10{]}）为：

\[
\frac {\partial \psi_ {k l}}{\partial a _ {m}} - \frac {\partial \psi_ {k m}}{\partial a _ {l}} = \frac {1}{2} \sum_ {1 \leqslant i, j \leqslant r} c _ {i j} ^ {k} (\psi_ {i l} \psi_ {j m} - \psi_ {j l} \psi_ {i m}) (1 \leqslant k, l, m \leqslant r).\tag{17}
\]

引入逆矩阵 \((\alpha_{ij})\)（它是 \((\psi_{ij})\) 的逆矩阵）以及无穷小变换

\[
\mathbf {X} _ {k} = \sum_ {i = 1} ^ {n} \xi_ {k i} \frac {\partial}{\partial x _ {i}}, \quad \mathbf {A} _ {k} = \sum_ {j = 1} ^ {r} \alpha_ {k j} \frac {\partial}{\partial a _ {j}} (1 \leqslant k \leqslant r)\tag{18}
\]

\begin{enumerate}
\def\labelenumi{(\arabic{enumi})}
\setcounter{enumi}{15}
\item
  并且 (17) 分别可写成：
\item
\end{enumerate}

\[
[ \mathrm{X} _ {i}, \mathrm{X} _ {j} ] = \sum_ {k = 1} ^ {r} c _ {i j} ^ {k} \mathrm{X} _ {k} (1 \leqslant i, j \leqslant r).\tag{20}
\]

\[
[ \mathbf {A} _ {i}, \mathbf {A} _ {j} ] = \sum_ {k = 1} ^ {r} c _ {i j} ^ {k} \mathbf {A} _ {k}
\]

反之，若给定 r 个无穷小变换 \(X_{k}\)（\(1 \leqslant k \leqslant r\)），它们线性无关且满足条件 (19)，则由这些变换生成的单参数子群生成一个具有 r 个本质参数的变换群。

最后，第三定理（{[}4{]}，第 1 卷，第 170、297 页以及第 3 卷，第 597 页）把满足 (19) 的无穷小变换系统 \((\mathbf{X}_k)_{1 \leqslant k \leqslant r}\) 的确定归结为一个纯代数问题：必须成立：

\begin{enumerate}
\def\labelenumi{(\arabic{enumi})}
\setcounter{enumi}{20}
\tightlist
\item
\end{enumerate}

\[
c _ {i j} ^ {k} + c _ {j i} ^ {k} = 0\tag{22}
\]

\[
\sum_ {l = 1} ^ {r} \left(c _ {i l} ^ {m} c _ {j k} ^ {l} + c _ {k i} ^ {m} c _ {i j} ^ {l} + c _ {j l} ^ {m} c _ {k l} ^ {l}\right) = 0 \quad (1 \leqslant i, j, k, m \leqslant r).
\]

反之，†若 (21) 和 (22) 成立，则存在一个满足关系 (19) 的无穷小变换系统，从而得到一个具有 r 个参数的变换群（换言之，\(X_{k}\) 的常系数组合构成一个李代数，反过来每个有限维李代数都可由此方式得到）。

这些结果还通过研究同构问题而得到补充。若能通过对变量作可逆坐标变换、对参数作可逆坐标变换而从一个变换群得到另一个，则称两个变换群相似；Lie 从研究一开始就在定义“典范参数”时自然地关注这一概念。他证明，若通过对“变量”作变换可以把一个群的无穷小变换化为另一个群的无穷小变换，则两个群相似（{[}4{]}，第 1 卷，第 329 页）。必要条件是这两个群的李代数同构，Lie 称这些群为“gleichzusammengesetzt”；但这一条件并不充分，因而有一整章（{[}4{]}，第 1 卷，第 19 章）专门讨论保证群“相似”的补充条件。另一方面，置换群理论提供了两个此类群“全同构”的概念（底层“抽象”群的同构）；Lie 将这一概念移植到变换群，并证明两个此类群“全同构”当且仅当它们的李代数同构（{[}4{]}，第 1 卷，第 418 页）。特别地，每个变换群都与它的每个参数群全同构；这说明，当我们希望研究群的结构时，它所作用于其上的“变量”并不重要，实际上重要的只有李代数。†

始终类比于置换群理论，Lie 引入了子群、正规子群和“满同态”（满射同态）的概念，并证明它们分别对应于子代数、理想和满射李代数同态；他此前已经遇到过一个特别重要的“满同态”例子，即伴随表示，并认识到它与群的中心的关系（{[}3f{]}，§ 3，第 9 号）。对于这些结果以及基本定理，关键工具是 Jacobi--Clebsch 定理，它给出了微分系统的完全可积性（这种定理的形式之一称为“{[}Frobenius 定理”）；不过，他使用单参数群给出了新的证明（{[}4{]}，第 1 卷，第 6 章）。

对置换群如此重要的传递性和本原性概念，同样自然地出现在“有限连续”变换群中，Lie--Engel 著作对此作了详细研究（{[}4{]}，第 1 卷，第 13 章及各处）；人们也认识到了它们与一点的稳定子群以及齐性空间概念之间的关系（在不采取全局观点所可能的范围内）（{[}4{]}，第 1 卷，第 425 页）。

最后，在 {[}4{]} 中，这部“词典”通过引入导出群和可解群的概念而得以完成（Lie 称之为“可积分群”；这一术语由微分方程理论启发，一直沿用到 H. Weyl 的著作）（{[}4{]}，第 1 卷，第 261 页及第 3 卷，第 678--679 页）；Lie 还在别处于 1883 年认识到了换位子与括号之间的关系（{[}3{]}，第 5 卷，第 358 页）。

\subsection*{基本定理的其他证明}

在 {[}8{]} 中，F. Schur 证明，在典范坐标中，(15) 的 \(\psi_{ik}\) 满足微分方程

\[
\frac {d}{d t} \left(t \psi_ {i k} (t a)\right) = \delta_ {i k} + \sum_ {j, l} c _ {j l} ^ {k} t a _ {l} \psi_ {i j} (t a).\tag{23}
\]

这些方程经积分后给出一个与下式等价的公式

\[
\varpi (\mathrm{X}) = \sum_ {n > 0} \frac {1}{(n + 1) !} (\mathrm{ad} (\mathrm{X})) ^ {n}\tag{24}
\]

本书第 III 章 §6 第 4 号命题 12 的公式；特别地，在典范坐标中，\(\psi_{ij}\) 可延拓为 \(a_k\) 的整函数。F. Schur 得出一个使 Lie 早先的评论更加精确的结果：若在变换群的定义 (4) 中只假定 \(f_i\) 属于 \(\mathbf{C}^2\) 类，则该群与解析群全同构。 \(\dagger\) E. Cartan 在研究微分系统积分的过程中，于 1904 年引入 Pfaff 形式（{[}12{]}，第 II 卷 \(_{2}\)，第 371 页）

\[
\omega_ {k} = \sum_ {i = 1} ^ {r} \psi_ {k i} d a _ {i} (1 \leqslant i \leqslant r)\tag{25}
\]

（在 (15) 的记号下），后来称为 Maurer--Cartan 形式。Maurer 条件 (17) 可写成

\[
d \omega_ {k} = - \frac {1}{2} \sum_ {i, j} c _ {i j} ^ {k} \omega_ {i} \wedge \omega_ {j};
\]

E. Cartan 证明，有限连续群理论可以从 \(\omega_{k}\) 出发建立，并确立了这一观点与 Lie 观点的等价性。但对他而言，这种方法的意义主要在于它可以适用于“无限连续群”，其理论被他推进得远超 Lie，并使他发展了广义“活动标架”理论。

\hnsection{李代数理论}

一旦建立了变换群与李代数之间的对应，理论便转向了更加代数化的方向，并以深入研究李代数为中心。†

1888 年至 1894 年是第一个短暂时期；这一时期以 Engel 及其学生 Umlauf，尤其是 Killing 和 E. Cartan 的工作为标志，在复李代数方面取得了一系列惊人的成果。我们此前已经看到，可解李代数的概念源于 Lie 本人；他曾证明（在复数情形下）可解线性李代数化为三角形形式的定理 ({[}4{]}, vol.~1, p.~270)。\(^{\ddagger}\) Killing 注意到 {[}11{]}，李代数中存在最大的可解理想（我们现在称之为根基），且该李代数对其根基的商的根基为零；他把根基为零的李代数称为半单李代数，并证明它们是单代数的乘积（后一个概念已经由 Lie 引入，他曾证明“经典”代数的单纯性 ({[}4{]}, vol.~3, p.~682)）。

另一方面，Killing 在李代数中引入了特征方程 \(\det(\mathrm{ad}(x)-\omega.1)=0\) ，Lie 在研究包含给定李代数元素的二维李子代数时已经遇到过这个方程。我们将在本书其他《历史注记》中回到对这些方法的分析：Killing 深入研究半单代数的“通用”特征方程的根的性质时，取得了他最卓越的成果，即完全确定（复）单李代数。§

Killing 证明，可解代数的导出代数是“秩为 0”的（这意味着对代数的每个元素 x，ad x 都是幂零的）。几乎与此同时，Engel 证明了“秩为 0”的代数是可解的（这个命题实质上就是我们在第一章 § 4，第 2 号中称为 Engel 定理的命题）。另一方面，E. Cartan 在其学位论文中引入了我们现在称为“Killing 型”的对象，并建立了两个基本判据，它们借助该型刻画了可解李代数和半单李代数。

Killing 曾断言 ({[}11{]}, IV)，李代数的导出代数是一个半单代数与其根基之和，而该根基是幂零的，但他的证明并不完整。不久以后，E. Cartan 未经证明便宣布 ({[}12{]}, vol.~I₁, p.~104)，更一般地说，每个李代数都是其根基与一个半单子代数之和；这一时期在这个方向上无可争议地确立的唯一结果，是 Engel 的一个定理，它断言每个非可解李代数中都存在一个 3 维单李子代数。Cartan 这一断言的第一个发表的证明（针对复李代数）归于 E. E. Levi {[}18{]}；J. H. C. Whitehead 于 1936 年给出了另一个证明（在实数情形下同样有效）{[}26 a{]}。1942 年，A. Malcev 通过关于“Levi 截面”在共轭意义下唯一性的定理完成了这一结果。

Lie 从其最初的著作起，就一直关注任意李代数与线性李代数之间的同构问题。他曾相信，通过考察伴随表示，他已经肯定地解决了这个问题（并由此推导出他的“第三定理”的一个证明）\([3\ d]\) ；后来他意识到自己的证明只对中心为零的李代数正确，又给出了另一个错误的证明（ \([3\ f]\) ，§ 3, no. 9），随后认识到（ \([3\ g]\) ，p.~231）这个问题仍未解决。事实上，这个问题长期悬而未决，直到 1935 年 Ado 才肯定地解决了它 {[}27{]}。另一方面，Lie 实质上提出了确定最小维数的单李代数的线性表示这一问题，并已对经典群解决了它；Cartan 在其学位论文中也解决了例外单代数的这一问题 \(^{\dagger}\) ；二十年后，他所采用的方法被他推广，用以得到实或复单李代数的所有不可约表示。

线性表示的完全可约性这一性质，似乎是 Study 首先（以几何形式）遇到的。在一篇未发表的手稿中（但已在 {[}4{]}，vol.~3, pp.~785--788 中引述），他证明了李代数 \(\mathbf{SL}(2,\mathbf{C})\) 的线性表示具有此性质，并对 \(\mathbf{SL}(3,\mathbf{C})\) 和 \(\mathbf{SL}(4,\mathbf{C})\) 取得了部分结果。Lie 和 Engel 借此猜想，\(\mathbf{SL}(n,\mathbf{C})\) 的完全可约性定理对所有 \(n\) 成立。半单李代数线性表示的完全可约性由 H. Weyl 于 1925 年† 通过一种全局型论证（见后文）建立。第一个代数证明由 Casimir 和 van der Waerden 于 1935 年取得 {[}32{]}；此后 R. Brauer {[}31{]}（这就是我们复现的证明）和 J. H. C. Whitehead {[}26 b{]} 又给出了其他代数证明。

最后，在研究指数映射的过程中（cf.~infra），H. Poincaré ({[}14{]}, vol.~3) 考察了由李代数各算子生成的全阶微分算子的结合代数；他实质上证明，若 \((\mathbf{X}_i)_{1\leqslant i\leqslant n}\) 是该李代数的一个基，则由 \(\mathbf{X}_i\) 生成的结合代数以 \(\mathbf{X}_i\) 的某些对称函数为基（这些函数是将给定单项式的因子作所有置换后得到的非交换“单项式”之和）。他的证明的本质是代数性的，这使他得以得到我们在第一章中抽象定义的包络代数结构。类似的证明由 G. Birkhoff {[}29 b{]} 和 E. Witt {[}30{]} 于 1937 年给出。

上述大多数结果限于实或复李代数，只有它们按通常意义对应于李群。Jacobson 开始研究定义在 \(\mathbf{R}\) 或 \(\mathbf{C}\) 以外的域上的李代数 {[}28 a{]}；他证明经典结果的大部分（即第一章的结果）对任意特征为零的域仍然成立。

\hnsection{指数映射与豪斯多夫公式}

关于指数映射的最初研究由 E. Study 和 F. Engel 完成；Engel {[}9 b{]} 注意到，对于 \(\mathbf{SL}(2, \mathbf{C})\)，指数映射不是满射（例如 \(\begin{pmatrix} -1 & a \\ 0 & -1 \end{pmatrix}\) 在 \(a \neq 0\) 时不是指数元），但对于 \(\mathbf{GL}(n, \mathbf{C})\) 是满射，因而对于 \(\mathbf{PGL}(n, \mathbf{C})\) 也是满射（后一个性质在 \(n = 2\) 时 Study 已经注意到）；由此，\(\mathbf{SL}(2, \mathbf{C})\) 和 \(\mathbf{PGL}(2, \mathbf{C})\) 给出了两个局部同构群的例子，然而从全局观点看它们却非常不同。Engel 还证明，其他经典群加上伸缩变换后指数映射也是满射；Maurer、Study 等人继续开展了这项工作，但没有产生实质性的新结果。

1899 年，H. Poincaré ({[}14{]}, vol.~3, pp.~169--172 and 173--212) 从另一个角度开始研究指数映射。他的论文似乎经过仓促编辑，因为他在数处断言连通群的每个元素都是指数元，而在别处却给出了相反的例子。他的结果主要涉及伴随表示：他证明，这样的群 G 的一个半单元素可能是李代数 L(G) 中无穷多个元素的指数，而一个非半单元素则可能根本不是指数元。如果 ad(X) 没有作为 2πi 的非零倍数的特征值，那么 exp 在 X 处是 étale 的。他还证明，若 U 和 V 在 L(G) 中描出回路，W 由连续性定义使得 \(e^{U} \cdot e^{V} = e^{W}\) ，也可能无法回到 W 的原始值。他使用了一个残数公式，其本质上相当于

\[
\Phi (\mathrm{adX}) = \frac {1}{2 \pi i} \int \frac {\Phi (\xi) d \xi}{\xi - \mathrm{adX}}
\]

其中 \(\operatorname{ad}(\mathbf{X})\) 是一个非零特征值重数均为 1 的半单元素，\(\Phi\) 是收敛半径足够大的整级数，积分沿包围 \(\operatorname{ad} X\) 的特征值的闭合曲线取；他还研究了当 X 趋近于具有重特征值的变换时所发生的情况。

关于在公式 \(e^{U}.e^{V}=e^{W}\) 中把 W 表示为 U 和 V 的函数的研究，在 Poincaré 的工作之前不久就已成为 Campbell 的两篇论文的主题 {[}13{]}。正如 Baker 不久后所写：“……李理论以显然的方式表明，乘积 \(e^{U}.e^{V}\) 的形式是 \(e^{W}\)，其中 W 是 U 和 V 的交错式级数……”。后来关于这一问题的工作旨在使这一断言精确化，并给出 W 的显式公式（或构造方法）（“豪斯多夫公式”）。继 Campbell 和 Poincaré 之后，Pascal、Baker {[}15{]} 和 Hausdorff {[}16{]} 又回到这个问题；他们各自都认为前人的证明并不令人信服；主要困难在于“交错式”究竟是什么意思：它们是所讨论李代数的特殊元素，还是普适的“符号”表达式？Campbell、Poincaré 和 Baker 都没有清楚地说明这一点。另一方面，Hausdorff 的论文极其精确；他首先研究有限个未定元的（非交换）结合形式幂级数代数，并把 U、V、W 看作这个代数的元素。他通过类似其前人的微分方程论证证明了 W 的存在。当未定元替换为有限维李代数的元素时，他又用同一论证证明了级数的收敛性。正如 Baker 所指出的（Poincaré 也独立地指出了这一点），这个结果可用于证明 Lie 的第三定理；他澄清了李群与李代数之间的对应关系，例如在换位子群方面。

1947 年，Dynkin {[}39{]} 再次研究了这个问题，并通过一开始就考察赋范李代数（无论有限维与否，定义在 R、C 或超度量域上），得到了豪斯多夫公式的显式系数。†

\hnsection{线性表示与全局李群}

刚才提到的这些工作没有任何一项真正试图定义和研究全局李群。H. Weyl 首先朝这个方向迈出了步伐。他受到两种理论的启发，这两种理论此前一直独立发展：E. Cartan 关于复半单李代数线性表示的理论，以及 Frobenius 关于有限群线性表示的理论；后者刚刚被 I. Schur 借助 Hurwitz 的一个想法转用于正交群。Hurwitz 已经说明 {[}17{]}，通过把有限群上的平均运算替换为关于不变测度的积分，可以为正交群或酉群构造不变量。他还注意到，将此方法应用于酉群，可以为一般线性群得到不变量，这就是“酉技巧”的第一个例子。1924 年，I. Schur {[}20{]} 用这一过程，通过构造不变的非退化正定 Hermite 型，证明了正交群 \(\mathbf{O}(n)\) 和酉群 \(\mathbf{U}(n)\) 的表示的完全可约性；他由“酉技巧”推出 \(\mathbf{O}(n,\mathbf{C})\) 和 \(\mathbf{SL}(n,\mathbf{C})\) 的全纯表示的完全可约性，建立了 \(\mathbf{O}(n)\) 和 \(\mathbf{U}(n)\) 的特征标正交关系，并确定了 \(\mathbf{O}(n)\) 的特征标。H. Weyl 随即将这种方法推广到复半单李代数 {[}21{]}。给定这样的代数 g，他证明它具有一个“紧实实形式”（这相当于说，它是由 R 上的一个代数 \(g_{0}\) 通过从 R 到 C 的标量扩张得到的，而其伴随群 \(G_{0}\) 是紧的）。此外，他证明 \(G_{0}\) 的基本群是有限的，因而 \(G_{0}\) 的万有覆盖‡是紧的。他通过适当改造 Schur 的方法，推出 g 的表示的完全可约性，并且还用一种全局方法确定了 g 的表示的特征标。在给 I. Schur 的一封信中（Sitzungsber, Berlin, 1924, pp.~338--343），H. Weyl 概述了 Cartan 的、Schur 尚不知道的结果（cf.~{[}20{]}，p.~299，页底注释），并比较了两种方法：Cartan 的方法给出了具有李代数 g 的单连通群的所有全纯表示；对于正交群，这样得到的是一个二重覆盖（后来称为自旋群）的表示，而 Schur 未能得到它；另一方面，Schur 的方法的优点在于显示了完全可约性并给出了特征标的显式形式。

在 H. Weyl 的工作之后，E. Cartan 在关于对称空间和李群的研究中采取了明确的全局观点。这一观点构成了他 1930 年对“有限连续”群理论的论述 ({[}12{]}, vol.~I₂, pp.~1165--1225) 的基础。特别是，这里给出了全局版本的第三基本定理（存在具有给定李代数的李群）的第一个证明；Cartan 还证明，实李群的每个闭子群都是李群（第三章，§ 8，第 2 号，定理 2），从而推广了 J. von Neumann 关于线性群的闭子群的一个结果 {[}23{]}。在这篇回忆录中，von Neumann 还证明了半单群的每个连续表示都是实解析的。

在这些工作之后，“经典”意义下的李群理论（即在 R 或 C 上的有限维理论）在其基本方面大体已经完成。Pontrjagin 在其关于拓扑群的著作中给出了第一个详细论述 {[}36{]}；在那里，他采用了与 Lie 很接近的观点，但仔细区分了局部与全局。随后是 Chevalley 的著作 {[}38{]}，其中还首次系统讨论了解析流形理论和外微分演算；Lie 的“无穷小变换”在那里作为向量场出现，李群 G 的李代数被认同为 G 上左不变向量场的空间。他不讨论“群芽”和“变换群”。

\hnsection{李群概念的扩展}

如今，李理论的活力体现在其应用的多样性（拓扑学、微分几何、算术等）以及平行理论的创立上；在这些理论中，底层流形结构被类似的结构所替代（ \(p\) 进流形、代数簇、概形、形式概形、\(\ldots\) ）。我们这里不讨论这些发展的历史，而只限于第三章涉及的内容：巴拿赫李群和 \(p\) 进李群。

\textbf{(a) 巴拿赫李群}

这里讨论的是“无限维”李群。从局部观点看，欧氏空间中的 0 的邻域被巴拿赫空间中的 0 的邻域所替代。这就是 G. Birkhoff 于 1936 年所做的工作 {[}29 a{]}，由此建立了完备赋范李代数的概念，以及它与定义在巴拿赫空间开集上的“群芽”之间的对应。约在 1950 年，Dynkin 将豪斯多夫公式推广到这一情形，从而完成了这些结果（cf.~supra）。

Birkhoff 和 Dynkin 的定义和结果都是局部的。直到最近，似乎还没有人试图明确建立相应的全局理论，这无疑是由于缺乏应用。†

\textbf{(b) p 进李群}

这类群最早见于 1907 年 Hensel 的工作 {[}19{]}，内容涉及 \(p\) 进解析函数（解析函数由整级数展开定义）。他特别研究了指数和对数；尽管定义它们的级数的先验行为令人惊讶（例如指数级数并不处处收敛），它们的基本函数性质仍然成立，从而在加法群与 \(\mathbf{Q}_p\) 的乘法群之间给出局部同构（或者更一般地，在任意特征为零的完备超度量域的加法群与乘法群之间给出局部同构）。

A. Weil {[}33{]} 和 E. Lutz {[}34{]} 在关于 p 进椭圆曲线（1936）的工作中同样研究了交换群（不过这次是非线性的）。除算术应用外，这里还基于对不变微分形式的积分，构造了群与加法群之间的局部同构。正如 C. Chabauty 不久后指出的，这一方法同样适用于阿贝尔簇；他未作进一步说明就用它证明了“莫德尔猜想”的一个特例 {[}35{]}。

从那时起，很明显，李群的局部理论几乎无需改变就可以应用于 p 进情形。李群--李代数“词典”的基本定理由 Chevalley 的学生 R. Hooke 在 1942 年的学位论文中建立 {[}37{]}；这项工作还包含 E. Cartan 关于实李群闭子群的定理的 p 进类比。

近来，M. Lazard {[}42 b{]} 为定义在 \(\mathbf{Q}_{\mathrm{p}}\) 上的紧解析群发展了更精确的“词典”形式。他证明，紧群 G 上存在 \(p\) 进解析结构，与 G 上某些滤过的存在密切相关，并由此给出各种应用（例如 G 的上同调）。Lazard 的工具之一，是改进 Dynkin 关于 \(p\) 进豪斯多夫级数收敛性的结果 {[}42 a{]}。

\hnsection{自由李代数}

最后还需谈及一系列关于李代数的工作，其中李代数与李群理论的联系十分微弱；另一方面，这些研究对于“抽象”群理论，尤其是幂零群理论，有着重要应用。

其源头是 P. Hall {[}24{]} 的工作，该工作发表于 1932 年。不过这里并未讨论李代数：P. Hall 想研究一类 p-群，即他称为“正则”的那些群。但这使他详细考察了群的迭代换位子和下中心列；关于此事，他建立了 Jacobi 恒等式的一个版本（cf.~第二章，§ 4, no. 4, 公式 (20)）以及“Hall 公式”

\[
(x y) ^ {n} = x ^ {n} y ^ {n} (x, y) ^ {n (1 - n) / 2} \dots \quad (\text { cf.   Chapter   II }, \S 5, \text { Exercise   9 }).
\]

几乎与此同时（1935--1937 年），W. Magnus ({[}25 a{]} 和 {[}25 b{]}) 以及 E. Witt {[}30{]} 的奠基性工作相继出现。在 {[}25 a{]} 中，Magnus 使用了与 Hausdorff 相同的形式幂级数代数 \(\hat{\mathbf{A}}\)（后来称为“Magnus 代数”）；他把自由群 F 嵌入其中，并利用 \(\hat{\mathbf{A}}\) 的自然滤过得到 F 的子群的一个递减序列 \((\mathbf{F}_{n})\)；这是滤过的最早例子之一。他猜想 \(F_{n}\) 与 F 的下中心列各项重合。这一猜想在他的第二篇论文中得到证明 {[}25 b{]}；在这项工作中，他还明确展示了自己的思想与 P. Hall 的思想之间的密切关系，并定义了自由李代数 L（作为 \(\hat{\mathbf{A}}\) 的子代数），他基本上证明了该代数可与 F 的分次代数认同。在 {[}30{]} 中，Witt 在多个方面完成了这一结果。特别是，他证明 L 的包络代数是自由结合代数，并立即推出 L 的齐次分量的秩（“Witt 公式”）。

至于 L 的基的证明，即所谓“Hall 基”（cf.~第二章，§ 2, no. 11），似乎直到 1950 年才出现在 M. Hall 的一篇短文中 {[}40{]}，尽管它已经隐含在上面引用的 P. Hall 和 W. Magnus 的工作中。

\specialchapter{书目}

\begin{enumerate}
\def\labelenumi{\arabic{enumi}.}
\item
  F. KLEIN and S. LIE: (a) Sur une certaine famille de courbes et surfaces, C. R. Acad. Sci., 70 (1870), pp.~1222--1226 and 1275--1279 (={[}2{]}, pp.~416--420 and {[}3{]}, vol.~1, pp.~78--85); (b) Über diejenigen ebenen Kurven, welche durch ein geschlossenes System von einfach unendlich vielen vertauschbaren linearen Transformationen in sich übergehen, Math. Ann., 4 (1871), pp.~50--84 (={[}2{]}, pp.~424--459 and {[}3{]}, vol.~1, Abh. XIV, pp.~229--266).
\item
  F. KLEIN, Gesammelte mathematische Abhandlungen, Bd. I, Berlin (Springer), 1921.
\item
  S. LIE, Gesammelte Abhandlungen, 7 vol., Leipzig (Teubner): (a) Über die Reziprozitätsverhältnisse des Reyeschen Komplexes, vol.~I, Abh. V, pp.~68--77 (=Gött. Nach. (1870), pp.~53--66); (b) Über eine Klasse geometrischer Transformationen, vol.~I, Abh. XII, pp.~153--214 (=Christiana For. (1871), pp.~182--245); (c) Über partielle Differentialgleichungen erster Ordnung, vol.~III, Abh. VII, pp.~32--63 (=Christiana For. (1873), pp.~16--51); (d) Theorie der Transformationsgruppen II, vol.~V, Abh. III, pp.~42--75 (=Archiv. Math., I (1876), pp.~152--193); (e) Ueber Gruppen von Transformationen, vol.~V, Abh. I, pp.~1--8 (=Gött. Nachr. (1874), pp.~529--542); (f) Allgemeine Untersuchungen über Differentialgleichungen, die eine kontinuierliche endliche Gruppe gestatten, vol.~VI, Abh. III, pp.~139--223 (=Math. Ann., 25 (1885), pp.~71--151); (g) Beiträge zur allgemeinen Transformationenstheorie, vol.~VI, Abh. V, pp.~230--236 (=Leipziger Ber. (1888), pp.~14--21); (h) Theorie der Transformationsgruppen III, vol.~V, Abh. IV, pp.~78--133 (=Archiv. Math., vol.~III, (1887), pp.~93--165).
\item
  S. LIE and F. ENGEL, Theorie der Transformationsgruppen, 3 vol., Leipzig (Teubner), 1888--1893.
\item
  J. J. SYLVESTER, Collected Mathematical Papers, 4 vols., Cambridge, 1904-1911.
\item
  A. CAYLEY, Collected Mathematical Papers, 13 vols., Cambridge, 1889-1898.
\item
  C. JORDAN, Mémoire sur les groupes de mouvements, Annali di Math., 11 (1868--1869), pp.~167--215 and 332--345 (= Oeuvres, vol.~IV, pp.~231--302).
\item
  F. SCHUR: (a) Zur Theorie der aus Haupteinheiten gebildeten Komplexen, Math. Ann., 33 (1889), pp.~49--60; (b) Neue Begründung der Theorie der endlichen Transformationsgruppen, Math. Ann., 35 (1890), pp.~161--197; (c) Zur Theorie der endlichen Transformationsgruppen, Math. Ann., 38 (1891), pp.~273--286; (d) Über den analytischen Character der eine endliche continuierliche Transformationsgruppe darstellende Funktionen, Math. Ann., 41 (1893), pp.~509--538.
\item
  F. ENGEL: (a) Über die Definitionsgleichung der continuerlichen Transformationsgruppen, Math. Ann., 27 (1886), pp.~1--57; (b) Die Erzeugung der endlichen Transformationen einer projektiven Gruppe durch die infinitesimalen Transformationen der Gruppe, I, Leipzig Ber., 44 (1892), pp.~279--296, II (mit Beiträgen von E. Study), ibid., 45 (1893), pp.~659--696.
\item
  L. MAURER, Über allgemeinere Invarianten-Systeme, Sitzungsber. München, 18 (1888), pp.~103--150.
\item
  W. KILLING, Die Zusammensetzung der stetigen endlichen Transformationsgruppen: (I) Math. Ann., 31 (1888), pp.~252--290; (II) ibid., 33 (1889), pp.~1--48; (III) ibid., 34 (1889), pp.~57--122; (IV) ibid., 36 (1890), pp.~161--189.
\item
  E. CARTAN, Oeuvres complètes, 6 vol., Paris (Gauthier-Villars), 1952--54.
\item
  J. E. CAMPBELL: (a) On a law of combination of operators bearing on the theory of continuous transformation groups, Proc. London Math. Soc., (1) 28 (1897), pp.~381--390; (b) On a law of combination of operators (second paper), ibid., 29 (1898), pp.~14--32.
\item
  H. POINCARÉ, Oeuvres, 11 vol., Paris (Gauthier-Villars), 1916--1956.
\item
  H. F. BAKER, Alternants and continuous groups, Proc. London Math. Soc., (2) 3 (1905), pp.~24--47.
\item
  F. HAUSDORFF, Die symbolische Exponentialformel in der Gruppentheorie, Leipziger Ber., 58 (1906), pp.~19--48.
\item
  A. HURWITZ, Über die Erzeugung der Invarianten durch Integration, Gött, Nachr. (1897), pp.~71--90 (=Math. Werke, vol.~II, pp.~546--564).
\item
  E. E. LEVI, Sulla struttura dei gruppi finiti e continui, Atti Acc. Sci. Torino, 40 (1905), pp.~551--565 (= Opere, vol.~I, pp.~101--115).
\item
  K. HENSEL, Über die arithmetischen Eigenschaften der Zahlen, Jahresber. der D.M.V., 16 (1907), pp.~299--319, 388--393, 474--496.
\item
  I. SCHUR, Neue Anwendungen der Integralrechnung auf Probleme der Invariantentheorie, Sitzungsber. Berlin, 1924, pp.~189--208, 297--321, 346--355.
\item
  H. WEYL, Theorie der Darstellung kontinuierlicher halb-einfacher Gruppen durch lineare Transformationen, I, Math. Zeitschr., 23, (1925),
\end{enumerate}

pp.~271--309; II, ibid., 24 (1926), pp.~328--376; III, ibid., 24 (1926), pp.~377--395 (=Werke, vol.~II, pp.~543--647).

\begin{enumerate}
\def\labelenumi{\arabic{enumi}.}
\setcounter{enumi}{21}
\item
  O. SCHREIER: (a) Abstrakte kontinuierliche Gruppen, Abh. math. Sem. Hamburg, 4 (1926), pp.~15--32; (b) Die Verwandschaft stetiger Gruppen in grossen, ibid., 5 (1927), pp.~233--244.
\item
  J. VON NEUMANN, Zur Theorie der Darstellung kontinuierlicher Gruppen, Sitzungsber. Berlin, 1927, pp.~76--90 (=Collected Works, vol.~I, pp.~134--148).
\item
  P. HALL, A contribution to the theory of groups of prime power order, Proc. London Math. Soc., (3) 4 (1932), pp.~29--95.
\item
  W. MAGNUS: (a) Beziehungen zwischen Gruppen und Idealen in einem speziellen Ring, Math. Ann., 111 (1935), pp.~259--280; (b) Über Beziehungen zwischen höheren Kommutatoren, J. Crelle, 177 (1937), pp.~105--115.
\item
  J. H. C. WHITEHEAD: (a) On the decomposition of an infinitesimal group, Proc. Camb. Phil. Soc., 32 (1936), pp.~229--237 (=Mathematical Works, I, pp.~281--289); (b) Certain equations in the algebra of a semi-simple infinitesimal group, Quart. Journ. of Math., (2) 8 (1937), pp.~220--237 (=Mathematical Works, I, pp.~291--308).
\item
  I. Ado: (a) Note on the representation of finite continuous groups by means of linear substitutions (俄文), Bull. Phys. Math. Soc. Kazan, 7 (1935), pp.~3--43; (b) The representation of Lie algebras by matrices (俄文), Uspehi Mat. Nauk, 2 (1947), pp.~159--173 (英文译本：Amer. Math. Soc. Transl., (1) 9, pp.~308--327).
\item
  N. JACOBSON: (a) Rational methods in the theory of Lie algebras, Ann. of Math., 36 (1935), pp.~875--881; (b) Classes of restricted Lie algebras of characteristic p, II, Duke Math. Journal, 10 (1943), pp.~107--121.
\item
  G. BIRKHOFF: (a) Continuous groups and linear spaces; Rec. Math. Moscou, 1 (1936), pp.~635--642; (b) Representability of Lie algebras and Lie groups by matrices, Ann. of Math., 38 (1937), pp.~526--532.
\item
  E. WITT, Treue Darstellung Lieschen Ringe, J. Creille, 177 (1937), pp.~152--160.
\item
  R. BRAUER, Eine Bedingung für vollständige Reduzibilität von Darstellungen gewöhnlicher und infinitesimaler Gruppen, Math. Zeitschr., 41 (1936), pp.~330--339.
\item
  H. CASIMIR-B. L. VAN DER WAERDEN, Algebraischer Beweis der vollständigen Reduzibilität der Darstellungen halbeinfacher Liescher Gruppen, Math. Ann., 111 (1935), pp.~1--12.
\item
  A. WEIL, Sur les fonctions elliptiques p-adiques, C. R. Acad. Sci., 203 (1935), p.~22.
\item
  E. Lutz, Sur l'équation \(y^{2} = x^{3} - Ax - B\) dans les corps \(p\) -adiques, \(J\) . Crelle, 177 (1937), pp.~237-247.
\item
  C. CHABAUTY, Sur les points rationels des courbes algébriques de genre supérieur à l'unité, C. R. Acad. Sci., 212 (1941), pp.~882--884.
\end{enumerate}

\begin{enumerate}
\def\labelenumi{\arabic{enumi}.}
\setcounter{enumi}{35}
\item
  L. S. PONTRJAGIN, Topological groups, Princeton Univ. Press, 1939.
\item
  R. HOOKE, Linear \(p\) -adic groups and their Lie algebras, Ann. of Math., 43 (1942), pp.~641-655.
\item
  C. CHEVALLEY, Theory of Lie groups, Princeton Univ. Press, 1946.
\item
  E. DYNKIN: (a) Evaluation of the coefficients of the Campbell-Hausdorff formula (俄文), Dokl. Akad. Nauk, 57 (1947), pp.~323-326; (b) Normed Lie algebras and analytic groups (俄文), Uspehi Mat. Nauk, 5 (1950), pp.~135-186 (英文译本：Amer. Math. Soc. Transl., (1) 9, pp.~470-534).
\item
  M. HALL, A basis for free Lie rings and higher commutators in free groups, Proc. Amer. Math. Soc., 1 (1950), pp.~575--581.
\item
  D. MONTGOMERY-L. ZIPPIN, Topological Transformation Groups, New York (Interscience), 1955.
\item
  M. LAZARD: (a) Quelques calculs concernant la formule de Haudorff, Bull. Soc. Math. France, 91 (1963), pp.~435--451; (b) Groupes analytiques p-adiques, Publ. Math. I.H.E.S., no. 26 (1965), pp.~389--603.
\end{enumerate}

\(\mathcal{C}_k\mathfrak{g}\)（g 为李代数）：I.1.6.

ad \(_{g}\) x, ad x（x 为李代数 g 的一个元素）：I.1.2.

\(g_{(K_{1})}\)（g 为李代数）：I.1.9.

\(e^u\) , \(\exp u\)（ \(u\) 为特征为 0 的域上向量空间的幂零自同态）：I.6.8.

\(\mathcal{D}\mathfrak{g}, \mathcal{D}^{\kappa}\mathfrak{g}, \mathcal{C}^{\kappa}\mathfrak{g}\)（g 为李代数）：I.1.5.
